% ================================================================
%  DOCUMENTAÇÃO COMPLETA — SISTEMA DE AGENDAMENTO
%  TCC — UNDB — Guilherme Braga
%
%  INSTRUÇÕES DE USO:
%  Cada bloco é precedido de comentário indicando o capítulo e a seção
%  de destino no documento principal.  Copie cada bloco para o
%  arquivo .tex correto ou inclua via % ================================================================
%  DOCUMENTAÇÃO COMPLETA — SISTEMA DE AGENDAMENTO
%  TCC — UNDB — Guilherme Braga
%
%  INSTRUÇÕES DE USO:
%  Cada bloco é precedido de comentário indicando o capítulo e a seção
%  de destino no documento principal.  Copie cada bloco para o
%  arquivo .tex correto ou inclua via % ================================================================
%  DOCUMENTAÇÃO COMPLETA — SISTEMA DE AGENDAMENTO
%  TCC — UNDB — Guilherme Braga
%
%  INSTRUÇÕES DE USO:
%  Cada bloco é precedido de comentário indicando o capítulo e a seção
%  de destino no documento principal.  Copie cada bloco para o
%  arquivo .tex correto ou inclua via % ================================================================
%  DOCUMENTAÇÃO COMPLETA — SISTEMA DE AGENDAMENTO
%  TCC — UNDB — Guilherme Braga
%
%  INSTRUÇÕES DE USO:
%  Cada bloco é precedido de comentário indicando o capítulo e a seção
%  de destino no documento principal.  Copie cada bloco para o
%  arquivo .tex correto ou inclua via \input{docs/tcc_completo}.
% ================================================================


% ################################################################
%  CAPÍTULO 1 — INTRODUÇÃO
%  (insira na seção de motivação / objetivos do capítulo 1)
% ################################################################

% ----------------------------------------------------------------
%  Parágrafo de contextualização — inserir após a apresentação do
%  problema de pesquisa
% ----------------------------------------------------------------
A interface de uma aplicação web é o primeiro ponto de contato entre
o usuário e o sistema; uma interface mal projetada pode comprometer a
adoção da ferramenta independentemente da qualidade do seu código de
negócio~\cite{nielsen1994}.  Nesse contexto, o presente trabalho dedica
atenção especial ao desenvolvimento de uma camada de apresentação
profissional, responsiva e acessível, utilizando tecnologias consolidadas
do ecossistema Java — Thymeleaf, Bootstrap 5 e Chart.js — sobre a
plataforma Spring Boot.

% ----------------------------------------------------------------
%  Objetivo específico — adicionar à lista de objetivos específicos
% ----------------------------------------------------------------
\begin{itemize}
  \item Projetar e implementar interfaces administrativas e um portal do
    cliente com padrões visuais modernos, responsividade \textit{mobile-first}
    e navegação intuitiva, adequadas à apresentação acadêmica e ao uso
    em ambiente de produção.
  \item Incorporar visualização de dados (\textit{dashboard} com gráfico de
    \textit{doughnut} e barras de progresso) que permita ao gestor monitorar
    a distribuição de agendamentos por status em tempo real.
\end{itemize}

% ----------------------------------------------------------------
%  Justificativa — inserir no parágrafo de justificativa
% ----------------------------------------------------------------
A escolha por \textit{server-side rendering} com Thymeleaf, em vez de um
\textit{framework} SPA (\textit{Single Page Application}), justifica-se pela
natureza do projeto: um sistema de agendamento com domínio bem definido,
cuja maioria das interações do usuário segue fluxos lineares (cadastro,
listagem, agendamento). Essa abordagem reduz a complexidade de implantação,
elimina a necessidade de uma API REST separada para as telas internas e
facilita a validação de formulários no servidor com o Spring MVC
\textit{BindingResult}.


% ################################################################
%  CAPÍTULO 2 — FUNDAMENTAÇÃO TEÓRICA
% ################################################################

% ----------------------------------------------------------------
%  Seção: Tecnologias de Backend
%  (inserir na subseção de Spring Boot / Spring MVC)
% ----------------------------------------------------------------
\subsection{Spring Boot e Spring MVC}

O Spring Boot~\cite{springboot} é um \textit{framework} que simplifica a
configuração e a implantação de aplicações Java baseadas no ecossistema
Spring.  Ele adota o paradigma de \textit{convention over configuration}:
por meio de dependências \textit{starter} e auto-configurações, o
desenvolvedor pode iniciar uma aplicação web produtiva com poucas linhas de
configuração.

O Spring MVC implementa o padrão arquitetural
\textit{Model-View-Controller}~\cite{gamma1994}: o \texttt{@Controller}
recebe a requisição HTTP, delega a lógica de negócio à camada de serviço e
devolve um nome de \textit{view} ao \texttt{DispatcherServlet}, que aciona
o motor de \textit{templates} Thymeleaf para gerar o HTML final enviado ao
navegador.

\subsection{Spring Security}

O Spring Security~\cite{springsecurity} fornece autenticação e autorização ao
sistema.  No contexto do painel administrativo, a autenticação é realizada
por meio de formulário HTML com sessão gerenciada pelo servidor; o
\texttt{SecurityConfig} configura as rotas protegidas (\texttt{/dashboard},
\texttt{/web/**}) e as rotas públicas (\texttt{/portal/**}).  O portal do
cliente utiliza autenticação manual via \texttt{HttpSession}, com senhas
armazenadas em \textit{hash} BCrypt (\texttt{BCryptPasswordEncoder}).

\subsection{Spring Data JPA e PostgreSQL}

A persistência é gerenciada pelo Spring Data JPA~\cite{springdatajpa},
que abstrai as operações de banco de dados por meio de interfaces
\texttt{JpaRepository}.  As entidades — \texttt{Cliente},
\texttt{Profissional}, \texttt{Servico} e \texttt{Agendamento} — são
mapeadas com anotações JPA (\texttt{@Entity}, \texttt{@ManyToOne},
\texttt{@Enumerated}) e o banco de dados utilizado é o PostgreSQL.

% ----------------------------------------------------------------
%  Seção: Tecnologias de Interface e Frontend
%  (inserir como subseção de Tecnologias Utilizadas)
% ----------------------------------------------------------------
\subsection{Tecnologias de Interface e \textit{Frontend}}

A camada de apresentação foi construída inteiramente sobre tecnologias
consolidadas do ecossistema Java, garantindo compatibilidade nativa com o
Spring Boot, manutenibilidade e facilidade de evolução futura.

\subsubsection{Thymeleaf}

O Thymeleaf~\cite{thymeleaf} é um motor de \textit{templates} para a
plataforma Java integrado nativamente ao Spring MVC.  Toda a geração de
HTML é realizada no servidor (\textit{server-side rendering}), eliminando a
necessidade de APIs REST separadas para as telas administrativas e o portal
do cliente.  Os \textit{templates} herdam de layouts-base por meio dos
mecanismos de fragmento (\texttt{th:replace} / \texttt{th:fragment}),
garantindo consistência visual e evitando duplicação de código.

Um aspecto técnico relevante explorado neste trabalho diz respeito à
interação entre fragmentos Thymeleaf e JavaScript: quando um fragmento é
passado como parâmetro a um \textit{layout} via
\texttt{th:replace="\textasciitilde\{layout :: html(\textasciitilde\{:: conteudo\})\}"},
qualquer diretiva \texttt{th:inline="javascript"} dentro desse fragmento
não é processada de forma confiável~\cite{thymeleaf}.  A solução adotada
foi injetar os dados do controlador como atributos \texttt{data-*} no
elemento HTML (\texttt{th:attr}) e lê-los no JavaScript puro por meio de
\texttt{element.dataset}, eliminando a dependência do processamento
\textit{inline} e mitigando riscos de XSS por interpolação direta.

\subsubsection{Bootstrap 5}

O Bootstrap 5.3.3~\cite{bootstrap5} fornece o sistema de \textit{grid}
responsivo baseado em \textit{Flexbox}, componentes visuais (Navbar, Modal,
Card, Badge, Table, Progress) e utilitários CSS utilizados em todas as
telas.  A versão 5 eliminou a dependência do jQuery, tornando o
\textit{framework} mais leve e compatível com o JavaScript puro utilizado
nos \textit{scripts} de validação e interação.

\subsubsection{Bootstrap Icons}

A biblioteca Bootstrap Icons 1.11.3~\cite{bootstrapicons} é empregada como
sistema de ícones SVG (\textit{icon font}), presente em botões, títulos de
página, itens de menu e \textit{badges} de status.  O uso consistente de
ícones contextuais acelera o reconhecimento de destinos de navegação e
reduz a carga cognitiva do usuário~\cite{pierce2004}.

\subsubsection{Chart.js}

O Chart.js 4.4.3~\cite{chartjs} é incorporado via CDN para renderizar o
gráfico de \textit{doughnut} de distribuição de agendamentos por status no
\textit{dashboard} administrativo.  Os dados são injetados pelo Thymeleaf
por meio de atributos \texttt{data-*} no elemento \texttt{<canvas>}
(Listagem~\ref{lst:chartjs-dataset}) e lidos no JavaScript por meio de
\texttt{canvas.dataset}, garantindo separação clara entre lógica de
servidor e apresentação no cliente.

\begin{lstlisting}[language=HTML, caption={Injeção de dados via atributos \texttt{data-*}
  e inicialização do Chart.js via \texttt{canvas.dataset}.},
  label=lst:chartjs-dataset, basicstyle=\small\ttfamily]
<!-- Template Thymeleaf: injeta dados no canvas -->
<canvas id="graficoAgendamentosStatus"
        th:attr="data-agendado=${totalAgendado},
                 data-confirmado=${totalConfirmado},
                 data-concluido=${totalConcluido},
                 data-cancelado=${totalCancelado}"
        style="max-height:220px;"></canvas>

<!-- JavaScript puro: le os dados sem th:inline -->
<script>
(function () {
    var c = document.getElementById('graficoAgendamentosStatus');
    var agendado   = parseInt(c.dataset.agendado   || '0', 10);
    var confirmado = parseInt(c.dataset.confirmado || '0', 10);
    var concluido  = parseInt(c.dataset.concluido  || '0', 10);
    var cancelado  = parseInt(c.dataset.cancelado  || '0', 10);
    new Chart(c, { type: 'doughnut',
        data: { labels: ['AGENDADO','CONFIRMADO','CONCLUIDO','CANCELADO'],
            datasets: [{ data: [agendado, confirmado, concluido, cancelado],
                backgroundColor: ['#667eea','#28a745','#6c757d','#dc3545'],
                borderWidth: 2, borderColor: '#fff' }] },
        options: { plugins: { legend: { display: false } },
            cutout: '65%', responsive: true,
            maintainAspectRatio: false } });
})();
</script>
\end{lstlisting}

% ----------------------------------------------------------------
%  Seção: Usabilidade e Design de Interface
% ----------------------------------------------------------------
\subsection{Usabilidade e \textit{Design} de Interface}

\subsubsection{Heurísticas de Nielsen}

Nielsen~\cite{nielsen1994} definiu dez heurísticas de usabilidade para o
\textit{design} de interfaces de usuário, amplamente adotadas como
referência na avaliação e no projeto de sistemas interativos.  Dentre as
heurísticas mais diretamente aplicadas neste trabalho destacam-se:
\textit{Consistency and Standards} (uso uniforme de cores, ícones e padrões
de interação), \textit{Recognition rather than recall} (ícones contextuais
nos menus) e \textit{User control and freedom} (botões de cancelar/voltar
em todos os formulários e modais de confirmação para ações destrutivas).

\subsubsection{Responsividade \textit{Mobile-First}}

A abordagem \textit{mobile-first} preconiza que o \textit{layout} base seja
projetado para telas pequenas e progressivamente aprimorado para telas
maiores por meio de \textit{media queries}~\cite{bootstrap5}.  O sistema
emprega o \textit{grid} Bootstrap (\texttt{col-sm-*}, \texttt{col-lg-*})
para adaptar todas as telas ao viewport de 375 px (smartphone) sem overflow
horizontal.


% ################################################################
%  CAPÍTULO 3 (ou 4) — METODOLOGIA
% ################################################################

% ----------------------------------------------------------------
%  Subseção: Implementação da Interface — inserir em Metodologia
% ----------------------------------------------------------------
\subsection{Implementação da Camada de Interface}

O desenvolvimento da camada de interface seguiu um processo iterativo
composto por quatro etapas principais:

\begin{enumerate}
  \item \textbf{Diagnóstico} — inspeção de cada tela existente identificando
    deficiências visuais, HTML malformado e ausência de padrão entre
    componentes.
  \item \textbf{Priorização} — as correções foram ordenadas por impacto no
    usuário: primeiro os defeitos estruturais (HTML inválido, rotas
    quebradas), em seguida melhorias de usabilidade (formulários, tabelas,
    modais) e por último aprimoramentos visuais e adição de funcionalidades
    novas (gráfico de \textit{doughnut}).
  \item \textbf{Implementação} — as alterações foram realizadas arquivo a
    arquivo, com compilação Maven (\texttt{mvn compile}) após cada grupo de
    mudanças para garantir ausência de erros de \textit{template}.
  \item \textbf{Verificação} — cada tela foi verificada segundo o
    \textit{checklist} de testes de interface (Apêndice~A), que cobre os
    fluxos principais e os estados de erro em resolução \textit{desktop}
    (1280 px) e \textit{mobile} (375 px).
\end{enumerate}

\subsection{Controle de Qualidade Estático}

O projeto é analisado automaticamente pelo Qodana~\cite{qodana}, ferramenta
de análise estática da JetBrains integrada ao pipeline CI/CD do repositório
GitHub.  Durante o desenvolvimento foram identificados e corrigidos dois
grupos de avisos:

\begin{itemize}
  \item \textbf{Bulk operation} — substituição de laços \texttt{for} que
    populavam mapas linha a linha por chamadas \texttt{stream().collect
    (Collectors.toMap(...))} e \texttt{stream().collect
    (Collectors.groupingBy(...))}, conforme recomendação de estilo da
    plataforma Java.
  \item \textbf{Nullability / data flow} — substituição de comparador
    lambda com acesso encadeado sem verificação de nulo por
    \texttt{Comparator.comparing(AgendamentoResponse::getDataHora)},
    garantindo que a referência de método é nula-segura em tempo de compilação.
\end{itemize}

\subsection{Estratégia de Ramificação (\textit{Branching})}

O desenvolvimento seguiu o fluxo GitHub Flow: cada conjunto de alterações é
realizado em uma \textit{branch} de funcionalidade (\texttt{feature/}),
submetida como \textit{Pull Request} para revisão automatizada (Qodana) e
mesclada à \textit{branch} principal (\texttt{main}) após aprovação dos
\textit{checks} de CI.


% ################################################################
%  CAPÍTULO DE RESULTADOS E DISCUSSÃO
% ################################################################

\section{Resultados e Discussão}

Esta seção apresenta os resultados obtidos com a implementação da camada de
interface, analisando cada módulo do sistema sob os aspectos funcional,
visual e de usabilidade.

% ......................
\subsection{Painel Administrativo}

\subsubsection{Login}

A tela de login (\texttt{/web/login}, Figura~\ref{fig:tela-login}) passou a
exibir um fundo escuro com gradiente e um \textit{card} centralizado de alta
profundidade visual (\texttt{box-shadow: 0 24px 60px rgba(0,0,0,0.4)}).  O
campo de senha conta com um botão de alternância de visibilidade
(\textit{password toggle}), funcionalidade recomendada pelo NIST SP
800-63B~\cite{nist80063b} para redução de erros de autenticação.

\begin{figure}[H]
  \centering
  \includegraphics[width=0.55\textwidth]{figuras/tela-login.png}
  \caption{Tela de login com gradiente escuro e alternância de senha.}
  \label{fig:tela-login}
\end{figure}

\subsubsection{Navbar e Layout Base}

A barra de navegação (\texttt{layout/base.html}, Figura~\ref{fig:admin-navbar})
foi reconstruída com Bootstrap Navbar responsiva, ícones semânticos
Bootstrap Icons em cada item de menu, destaque dinâmico do item ativo por
JavaScript e colapso em \textit{hamburger menu} em telas menores que 992 px.
Alertas de sucesso e erro são exibidos com ícone e botão de fechamento
(\texttt{data-bs-dismiss="alert"}).

\begin{figure}[H]
  \centering
  \includegraphics[width=0.95\textwidth]{figuras/admin-navbar.png}
  \caption{Navbar administrativa com ícones e responsividade.}
  \label{fig:admin-navbar}
\end{figure}

\subsubsection{\textit{Dashboard}}

O \textit{dashboard} (\texttt{/dashboard}, Figura~\ref{fig:dashboard})
concentra as métricas operacionais mais relevantes em uma única tela:

\begin{itemize}
  \item \textbf{Quatro cartões de métricas} com fundo em gradiente, valor
    numérico em tipografia \textit{bold} de grande porte e efeito
    \textit{hover} de elevação (transform translateY(-4px)).
  \item \textbf{Gráfico de \textit{doughnut}} (Chart.js) com distribuição
    percentual dos agendamentos por status (Figura~\ref{fig:dashboard-chart}).
  \item \textbf{Barras de progresso Bootstrap} como visualização alternativa
    nativa do servidor (independente de JavaScript), exibindo contagem e
    percentual de cada status.
  \item \textbf{Tabela de próximos agendamentos} — até três registros com
    status AGENDADO ou CONFIRMADO, ordenados cronologicamente.
  \item \textbf{Acesso rápido} — três \textit{cards} com contadores de
    clientes, profissionais e serviços, com botões \textit{Gerenciar} e
    \textit{Novo} diretos.
\end{itemize}

\begin{figure}[H]
  \centering
  \includegraphics[width=0.95\textwidth]{figuras/dashboard.png}
  \caption{\textit{Dashboard} com métricas, gráfico e acesso rápido.}
  \label{fig:dashboard}
\end{figure}

\begin{figure}[H]
  \centering
  \includegraphics[width=0.55\textwidth]{figuras/dashboard-chart.png}
  \caption{Gráfico \textit{doughnut} de distribuição de agendamentos por status.}
  \label{fig:dashboard-chart}
\end{figure}

A Tabela~\ref{tab:cores-status} detalha a paleta de cores adotada para os
status, consistente entre o gráfico, as barras de progresso e os
\textit{badges} das tabelas.

\begin{table}[H]
\centering
\begin{tabular}{lll}
\hline
\textbf{Status} & \textbf{Cor}  & \textbf{Código Hexadecimal} \\ \hline
AGENDADO        & Azul-violeta  & \texttt{\#667eea}           \\
CONFIRMADO      & Verde         & \texttt{\#28a745}           \\
CONCLUÍDO       & Cinza         & \texttt{\#6c757d}           \\
CANCELADO       & Vermelho      & \texttt{\#dc3545}           \\ \hline
\end{tabular}
\caption{Paleta de cores dos status de agendamento.}
\label{tab:cores-status}
\end{table}

\paragraph{Desafio técnico — Chart.js em fragmento Thymeleaf.}
A integração do Chart.js com Thymeleaf exigiu a resolução de um problema
não documentado de forma proeminente: a diretiva \texttt{th:inline="javascript"}
não é processada quando o fragmento que a contém é passado como parâmetro a
um \textit{layout} template.  A solução implementada utilizou atributos
\texttt{data-*} no elemento \texttt{<canvas>} (injetados via \texttt{th:attr}
no servidor) e leitura via \texttt{canvas.dataset} no JavaScript do cliente,
conforme descrito na Listagem~\ref{lst:chartjs-dataset}.  Esta abordagem é
igualmente recomendada por questões de segurança: evita interpolação direta
de dados em contexto JavaScript, que poderia introduzir vulnerabilidades XSS.

\subsubsection{Listagens Administrativas}

Todas as telas de listagem seguem um padrão visual unificado
(Figuras~\ref{fig:clientes-lista},~\ref{fig:agendamentos-lista}):
cabeçalho com contador e botão de criação; busca em tempo real via JavaScript
(clientes, profissionais, serviços) ou filtro por formulário GET (agendamentos);
tabela Bootstrap \texttt{table-hover} com botões de ação ícone-apenas;
modal de confirmação para exclusão.

A listagem de agendamentos apresentava HTML malformado na versão anterior —
\texttt{<table>} sem classes Bootstrap, \texttt{</div>} órfãos e referências
a modais inexistentes.  A tela foi reescrita na íntegra, preservando toda a
lógica de transição de status e exclusão.

\begin{figure}[H]
  \centering
  \includegraphics[width=0.95\textwidth]{figuras/clientes-lista.png}
  \caption{Listagem de clientes com busca em tempo real e ações por ícone.}
  \label{fig:clientes-lista}
\end{figure}

\begin{figure}[H]
  \centering
  \includegraphics[width=0.95\textwidth]{figuras/agendamentos-lista.png}
  \caption{Listagem de agendamentos com filtros, badges de status e modais.}
  \label{fig:agendamentos-lista}
\end{figure}

\subsection{Portal do Cliente}

\subsubsection{\textit{Home} do Portal}

A página inicial do portal (\texttt{/portal}, Figura~\ref{fig:portal-home})
foi reformulada com um cartão \textit{hero} em gradiente
(\texttt{linear-gradient(135deg, \#667eea 0\%, \#764ba2 100\%)}), substituindo
o \textit{layout} simples anterior.  Três \textit{cards} de estatísticas
exibem serviços disponíveis, total de profissionais e agendamentos do cliente
autenticado.

\begin{figure}[H]
  \centering
  \includegraphics[width=0.95\textwidth]{figuras/portal-home.png}
  \caption{\textit{Home} do portal com \textit{hero card} e estatísticas.}
  \label{fig:portal-home}
\end{figure}

\subsubsection{\textit{Wizard} de Agendamento}

O fluxo de agendamento em quatro etapas
(Figuras~\ref{fig:wizard-step1},~\ref{fig:wizard-step3},~\ref{fig:wizard-step4})
guia o cliente desde a seleção do serviço até a confirmação com número de
protocolo.  O indicador de progresso visual (\textit{stepper}) comunica em
qual etapa o cliente se encontra e quais já foram concluídas (ícone de
verificação verde).  A Etapa 1 continha dois blocos \texttt{<script>}
duplicados que registravam os mesmos \textit{event listeners}; o bloco
redundante foi removido, eliminando a execução duplicada das funções de
seleção.

\begin{figure}[H]
  \centering
  \includegraphics[width=0.75\textwidth]{figuras/wizard-step1.png}
  \caption{Etapa 1 do \textit{wizard}: seleção de serviço por \textit{card}.}
  \label{fig:wizard-step1}
\end{figure}

\begin{figure}[H]
  \centering
  \includegraphics[width=0.65\textwidth]{figuras/wizard-step3.png}
  \caption{Etapa 3 do \textit{wizard}: revisão e confirmação do agendamento.}
  \label{fig:wizard-step3}
\end{figure}

\begin{figure}[H]
  \centering
  \includegraphics[width=0.65\textwidth]{figuras/wizard-step4.png}
  \caption{Etapa 4: tela de sucesso com protocolo de confirmação.}
  \label{fig:wizard-step4}
\end{figure}

\subsection{Melhorias de Usabilidade — Síntese}

A Tabela~\ref{tab:melhorias} resume as principais melhorias implementadas,
associando cada uma à heurística de Nielsen correspondente.

\begin{table}[H]
\centering
\small
\begin{tabular}{p{5cm} p{4cm} p{4.5cm}}
\hline
\textbf{Melhoria} & \textbf{Heurística} & \textbf{Resultado} \\ \hline
Fundo gradiente + card sombra na tela de login &
  Aesthetic \& minimalist design &
  Foco no formulário, redução de distração \\[4pt]
Ícones nos itens de menu &
  Recognition rather than recall &
  Navegação mais rápida \\[4pt]
\textit{Password toggle} &
  Error prevention &
  Menos erros de digitação de senha \\[4pt]
Modal de confirmação para exclusão &
  User control and freedom &
  Reversibilidade de ações destrutivas \\[4pt]
Gráfico + barras de progresso no dashboard &
  Visibility of system status &
  Distribuição de agendamentos em tempo real \\[4pt]
Busca em tempo real nas listagens &
  Efficiency of use &
  Localização rápida de registros \\[4pt]
Correção HTML malformado (agendamentos) &
  Consistency and standards &
  Comportamento previsível dos componentes \\ \hline
\end{tabular}
\caption{Melhorias de usabilidade e heurísticas de Nielsen correspondentes.}
\label{tab:melhorias}
\end{table}

\subsection{Discussão}

Os resultados demonstram que a adoção de um \textit{design system} baseado
em variáveis CSS globais e nos componentes Bootstrap 5 foi fundamental para
manter a consistência visual entre as mais de quinze telas do sistema, sem
impacto na lógica de negócio.  A separação entre camada de apresentação
(Thymeleaf + Bootstrap) e camada de negócio (serviços e repositórios Spring)
permaneceu preservada em todas as alterações realizadas.

A decisão de oferecer barras de progresso Bootstrap como alternativa nativa
ao gráfico Chart.js mostrou-se acertada sob dois aspectos: primeiro, garante
que o \textit{dashboard} exiba informações de status mesmo em ambientes com
JavaScript desabilitado ou bloqueado; segundo, evidencia os dados numéricos
de forma inequívoca, complementando a visualização proporcional do gráfico.

O desafio técnico da integração Chart.js/Thymeleaf ilustra um aspecto
relevante do desenvolvimento com \textit{server-side rendering}: a
necessidade de compreender os limites do processamento de \textit{templates}
no lado do servidor antes de tomar decisões de integração com bibliotecas
JavaScript.  A solução por atributos \texttt{data-*} é idiomática na web
moderna e agnóstica ao \textit{framework} de \textit{templates}, podendo ser
reaproveitada em projetos futuros que utilizem Freemarker, Pebble ou JSP.


% ################################################################
%  CAPÍTULO — CONCLUSÃO
% ################################################################

\section{Conclusão}

Este trabalho apresentou o desenvolvimento de um sistema de agendamento
construído sobre a plataforma Spring Boot, com ênfase na qualidade da camada
de interface e na experiência do usuário.  O sistema atende ao objetivo
principal de oferecer um painel administrativo completo e um portal do
cliente com fluxo de agendamento guiado, ambos responsivos e visualmente
adequados à apresentação acadêmica e ao ambiente de produção.

\subsection{Contribuições}

As principais contribuições técnicas do trabalho são:

\begin{itemize}
  \item \textbf{Arquitetura de \textit{templates} Thymeleaf} com fragmentos e
    layouts-base reutilizáveis, demonstrando boas práticas de organização de
    \textit{templates} em projetos Spring Boot.
  \item \textbf{Padrão de integração Chart.js/Thymeleaf} via atributos
    \texttt{data-*}, documentando e resolvendo a incompatibilidade entre
    \texttt{th:inline="javascript"} e fragmentos parametrizados, que não
    consta de forma destacada na documentação oficial do Thymeleaf.
  \item \textbf{\textit{Design system} coerente} — paleta, espaçamentos, raios
    de borda e sombras definidos em variáveis CSS (\texttt{:root}), facilmente
    adaptável a outras identidades visuais.
  \item \textbf{Fluxo de agendamento em \textit{wizard}} com indicador de
    progresso, validação por etapa e tela de confirmação com protocolo,
    padrão aplicável a outros domínios de negócio.
  \item \textbf{Pipeline de qualidade estática} — correção de avisos Qodana
    sobre operações de coleção e nulabilidade, demonstrando integração de
    análise estática no fluxo de desenvolvimento.
\end{itemize}

\subsection{Limitações}

O trabalho apresenta as seguintes limitações que podem ser endereçadas em
trabalhos futuros:

\begin{itemize}
  \item \textbf{Testes de interface automatizados} — o \textit{checklist} de
    testes foi realizado manualmente.  A automação com Selenium ou Playwright
    aumentaria a cobertura e a rastreabilidade das verificações.
  \item \textbf{Cobertura de testes de unidade} — os testes automatizados
    existentes cobrem os serviços principais; a camada de \textit{controllers}
    não possui testes de integração com MockMvc.
  \item \textbf{Acessibilidade formal} — embora as interfaces sigam
    recomendações gerais de contraste WCAG 2.1, não foi realizada auditoria
    formal com ferramentas como Axe ou Lighthouse.
  \item \textbf{Internacionalização} — o sistema está fixado em português
    brasileiro; suporte a múltiplos idiomas exigiria a adoção do mecanismo
    \texttt{MessageSource} do Spring.
\end{itemize}

\subsection{Trabalhos Futuros}

Como extensões naturais do sistema, destacam-se:

\begin{itemize}
  \item Implementação de notificações por e-mail e/ou WhatsApp para lembrete
    de agendamentos com proximidade da data.
  \item Exposição de uma API REST com autenticação JWT para integração com
    aplicativos móveis nativos.
  \item Adição de relatórios exportáveis (PDF/CSV) de agendamentos por
    período, profissional ou serviço.
  \item Suporte a múltiplos estabelecimentos (multi-tenant) na mesma
    instalação.
\end{itemize}

Em síntese, o sistema de agendamento desenvolvido neste trabalho cumpre seu
objetivo acadêmico e funcional, entregando uma plataforma operacional com
interfaces de qualidade profissional, arquitetura limpa e pipeline de
qualidade integrado.


% ################################################################
%  REFERÊNCIAS — adicionar ao arquivo .bib do projeto
% ################################################################

%
% @misc{springboot,
%   author = {{Pivotal Software / VMware}},
%   title  = {Spring Boot Reference Documentation},
%   year   = {2024},
%   url    = {https://docs.spring.io/spring-boot/docs/current/reference/html}
% }
% @misc{springsecurity,
%   author = {{Spring Security Team}},
%   title  = {Spring Security Reference},
%   year   = {2024},
%   url    = {https://docs.spring.io/spring-security/reference}
% }
% @misc{springdatajpa,
%   author = {{Spring Data Team}},
%   title  = {Spring Data JPA Reference},
%   year   = {2024},
%   url    = {https://docs.spring.io/spring-data/jpa/docs/current/reference/html}
% }
% @misc{thymeleaf,
%   author = {{Thymeleaf Team}},
%   title  = {Thymeleaf 3.x Documentation},
%   year   = {2024},
%   url    = {https://www.thymeleaf.org/documentation.html}
% }
% @misc{bootstrap5,
%   author = {{The Bootstrap Authors}},
%   title  = {Bootstrap 5.3 Documentation},
%   year   = {2024},
%   url    = {https://getbootstrap.com/docs/5.3}
% }
% @misc{bootstrapicons,
%   author = {{The Bootstrap Authors}},
%   title  = {Bootstrap Icons 1.11},
%   year   = {2024},
%   url    = {https://icons.getbootstrap.com}
% }
% @misc{chartjs,
%   author = {{Chart.js Contributors}},
%   title  = {Chart.js 4.x Documentation},
%   year   = {2024},
%   url    = {https://www.chartjs.org/docs/latest}
% }
% @misc{qodana,
%   author = {{JetBrains}},
%   title  = {Qodana — Static Code Analysis},
%   year   = {2024},
%   url    = {https://www.jetbrains.com/qodana}
% }
% @book{nielsen1994,
%   author    = {Nielsen, Jakob},
%   title     = {Usability Engineering},
%   publisher = {Morgan Kaufmann},
%   year      = {1994}
% }
% @techreport{wcag21,
%   author      = {{W3C Web Content Accessibility Guidelines Working Group}},
%   title       = {{WCAG} 2.1},
%   institution = {W3C},
%   year        = {2018},
%   url         = {https://www.w3.org/TR/WCAG21}
% }
% @techreport{nist80063b,
%   author      = {Grassi, Paul A. and others},
%   title       = {{NIST SP 800-63B}: Digital Identity Guidelines -- Authentication},
%   institution = {NIST},
%   year        = {2017}
% }
% @article{pierce2004,
%   author  = {Pierce, Jon L.},
%   title   = {The Role of Icons in User Interface Design},
%   journal = {Journal of Usability Studies},
%   year    = {2004}
% }
% @book{gamma1994,
%   author    = {Gamma, Erich and Helm, Richard and Johnson, Ralph and Vlissides, John},
%   title     = {Design Patterns: Elements of Reusable Object-Oriented Software},
%   publisher = {Addison-Wesley},
%   year      = {1994}
% }
.
% ================================================================


% ################################################################
%  CAPÍTULO 1 — INTRODUÇÃO
%  (insira na seção de motivação / objetivos do capítulo 1)
% ################################################################

% ----------------------------------------------------------------
%  Parágrafo de contextualização — inserir após a apresentação do
%  problema de pesquisa
% ----------------------------------------------------------------
A interface de uma aplicação web é o primeiro ponto de contato entre
o usuário e o sistema; uma interface mal projetada pode comprometer a
adoção da ferramenta independentemente da qualidade do seu código de
negócio~\cite{nielsen1994}.  Nesse contexto, o presente trabalho dedica
atenção especial ao desenvolvimento de uma camada de apresentação
profissional, responsiva e acessível, utilizando tecnologias consolidadas
do ecossistema Java — Thymeleaf, Bootstrap 5 e Chart.js — sobre a
plataforma Spring Boot.

% ----------------------------------------------------------------
%  Objetivo específico — adicionar à lista de objetivos específicos
% ----------------------------------------------------------------
\begin{itemize}
  \item Projetar e implementar interfaces administrativas e um portal do
    cliente com padrões visuais modernos, responsividade \textit{mobile-first}
    e navegação intuitiva, adequadas à apresentação acadêmica e ao uso
    em ambiente de produção.
  \item Incorporar visualização de dados (\textit{dashboard} com gráfico de
    \textit{doughnut} e barras de progresso) que permita ao gestor monitorar
    a distribuição de agendamentos por status em tempo real.
\end{itemize}

% ----------------------------------------------------------------
%  Justificativa — inserir no parágrafo de justificativa
% ----------------------------------------------------------------
A escolha por \textit{server-side rendering} com Thymeleaf, em vez de um
\textit{framework} SPA (\textit{Single Page Application}), justifica-se pela
natureza do projeto: um sistema de agendamento com domínio bem definido,
cuja maioria das interações do usuário segue fluxos lineares (cadastro,
listagem, agendamento). Essa abordagem reduz a complexidade de implantação,
elimina a necessidade de uma API REST separada para as telas internas e
facilita a validação de formulários no servidor com o Spring MVC
\textit{BindingResult}.


% ################################################################
%  CAPÍTULO 2 — FUNDAMENTAÇÃO TEÓRICA
% ################################################################

% ----------------------------------------------------------------
%  Seção: Tecnologias de Backend
%  (inserir na subseção de Spring Boot / Spring MVC)
% ----------------------------------------------------------------
\subsection{Spring Boot e Spring MVC}

O Spring Boot~\cite{springboot} é um \textit{framework} que simplifica a
configuração e a implantação de aplicações Java baseadas no ecossistema
Spring.  Ele adota o paradigma de \textit{convention over configuration}:
por meio de dependências \textit{starter} e auto-configurações, o
desenvolvedor pode iniciar uma aplicação web produtiva com poucas linhas de
configuração.

O Spring MVC implementa o padrão arquitetural
\textit{Model-View-Controller}~\cite{gamma1994}: o \texttt{@Controller}
recebe a requisição HTTP, delega a lógica de negócio à camada de serviço e
devolve um nome de \textit{view} ao \texttt{DispatcherServlet}, que aciona
o motor de \textit{templates} Thymeleaf para gerar o HTML final enviado ao
navegador.

\subsection{Spring Security}

O Spring Security~\cite{springsecurity} fornece autenticação e autorização ao
sistema.  No contexto do painel administrativo, a autenticação é realizada
por meio de formulário HTML com sessão gerenciada pelo servidor; o
\texttt{SecurityConfig} configura as rotas protegidas (\texttt{/dashboard},
\texttt{/web/**}) e as rotas públicas (\texttt{/portal/**}).  O portal do
cliente utiliza autenticação manual via \texttt{HttpSession}, com senhas
armazenadas em \textit{hash} BCrypt (\texttt{BCryptPasswordEncoder}).

\subsection{Spring Data JPA e PostgreSQL}

A persistência é gerenciada pelo Spring Data JPA~\cite{springdatajpa},
que abstrai as operações de banco de dados por meio de interfaces
\texttt{JpaRepository}.  As entidades — \texttt{Cliente},
\texttt{Profissional}, \texttt{Servico} e \texttt{Agendamento} — são
mapeadas com anotações JPA (\texttt{@Entity}, \texttt{@ManyToOne},
\texttt{@Enumerated}) e o banco de dados utilizado é o PostgreSQL.

% ----------------------------------------------------------------
%  Seção: Tecnologias de Interface e Frontend
%  (inserir como subseção de Tecnologias Utilizadas)
% ----------------------------------------------------------------
\subsection{Tecnologias de Interface e \textit{Frontend}}

A camada de apresentação foi construída inteiramente sobre tecnologias
consolidadas do ecossistema Java, garantindo compatibilidade nativa com o
Spring Boot, manutenibilidade e facilidade de evolução futura.

\subsubsection{Thymeleaf}

O Thymeleaf~\cite{thymeleaf} é um motor de \textit{templates} para a
plataforma Java integrado nativamente ao Spring MVC.  Toda a geração de
HTML é realizada no servidor (\textit{server-side rendering}), eliminando a
necessidade de APIs REST separadas para as telas administrativas e o portal
do cliente.  Os \textit{templates} herdam de layouts-base por meio dos
mecanismos de fragmento (\texttt{th:replace} / \texttt{th:fragment}),
garantindo consistência visual e evitando duplicação de código.

Um aspecto técnico relevante explorado neste trabalho diz respeito à
interação entre fragmentos Thymeleaf e JavaScript: quando um fragmento é
passado como parâmetro a um \textit{layout} via
\texttt{th:replace="\textasciitilde\{layout :: html(\textasciitilde\{:: conteudo\})\}"},
qualquer diretiva \texttt{th:inline="javascript"} dentro desse fragmento
não é processada de forma confiável~\cite{thymeleaf}.  A solução adotada
foi injetar os dados do controlador como atributos \texttt{data-*} no
elemento HTML (\texttt{th:attr}) e lê-los no JavaScript puro por meio de
\texttt{element.dataset}, eliminando a dependência do processamento
\textit{inline} e mitigando riscos de XSS por interpolação direta.

\subsubsection{Bootstrap 5}

O Bootstrap 5.3.3~\cite{bootstrap5} fornece o sistema de \textit{grid}
responsivo baseado em \textit{Flexbox}, componentes visuais (Navbar, Modal,
Card, Badge, Table, Progress) e utilitários CSS utilizados em todas as
telas.  A versão 5 eliminou a dependência do jQuery, tornando o
\textit{framework} mais leve e compatível com o JavaScript puro utilizado
nos \textit{scripts} de validação e interação.

\subsubsection{Bootstrap Icons}

A biblioteca Bootstrap Icons 1.11.3~\cite{bootstrapicons} é empregada como
sistema de ícones SVG (\textit{icon font}), presente em botões, títulos de
página, itens de menu e \textit{badges} de status.  O uso consistente de
ícones contextuais acelera o reconhecimento de destinos de navegação e
reduz a carga cognitiva do usuário~\cite{pierce2004}.

\subsubsection{Chart.js}

O Chart.js 4.4.3~\cite{chartjs} é incorporado via CDN para renderizar o
gráfico de \textit{doughnut} de distribuição de agendamentos por status no
\textit{dashboard} administrativo.  Os dados são injetados pelo Thymeleaf
por meio de atributos \texttt{data-*} no elemento \texttt{<canvas>}
(Listagem~\ref{lst:chartjs-dataset}) e lidos no JavaScript por meio de
\texttt{canvas.dataset}, garantindo separação clara entre lógica de
servidor e apresentação no cliente.

\begin{lstlisting}[language=HTML, caption={Injeção de dados via atributos \texttt{data-*}
  e inicialização do Chart.js via \texttt{canvas.dataset}.},
  label=lst:chartjs-dataset, basicstyle=\small\ttfamily]
<!-- Template Thymeleaf: injeta dados no canvas -->
<canvas id="graficoAgendamentosStatus"
        th:attr="data-agendado=${totalAgendado},
                 data-confirmado=${totalConfirmado},
                 data-concluido=${totalConcluido},
                 data-cancelado=${totalCancelado}"
        style="max-height:220px;"></canvas>

<!-- JavaScript puro: le os dados sem th:inline -->
<script>
(function () {
    var c = document.getElementById('graficoAgendamentosStatus');
    var agendado   = parseInt(c.dataset.agendado   || '0', 10);
    var confirmado = parseInt(c.dataset.confirmado || '0', 10);
    var concluido  = parseInt(c.dataset.concluido  || '0', 10);
    var cancelado  = parseInt(c.dataset.cancelado  || '0', 10);
    new Chart(c, { type: 'doughnut',
        data: { labels: ['AGENDADO','CONFIRMADO','CONCLUIDO','CANCELADO'],
            datasets: [{ data: [agendado, confirmado, concluido, cancelado],
                backgroundColor: ['#667eea','#28a745','#6c757d','#dc3545'],
                borderWidth: 2, borderColor: '#fff' }] },
        options: { plugins: { legend: { display: false } },
            cutout: '65%', responsive: true,
            maintainAspectRatio: false } });
})();
</script>
\end{lstlisting}

% ----------------------------------------------------------------
%  Seção: Usabilidade e Design de Interface
% ----------------------------------------------------------------
\subsection{Usabilidade e \textit{Design} de Interface}

\subsubsection{Heurísticas de Nielsen}

Nielsen~\cite{nielsen1994} definiu dez heurísticas de usabilidade para o
\textit{design} de interfaces de usuário, amplamente adotadas como
referência na avaliação e no projeto de sistemas interativos.  Dentre as
heurísticas mais diretamente aplicadas neste trabalho destacam-se:
\textit{Consistency and Standards} (uso uniforme de cores, ícones e padrões
de interação), \textit{Recognition rather than recall} (ícones contextuais
nos menus) e \textit{User control and freedom} (botões de cancelar/voltar
em todos os formulários e modais de confirmação para ações destrutivas).

\subsubsection{Responsividade \textit{Mobile-First}}

A abordagem \textit{mobile-first} preconiza que o \textit{layout} base seja
projetado para telas pequenas e progressivamente aprimorado para telas
maiores por meio de \textit{media queries}~\cite{bootstrap5}.  O sistema
emprega o \textit{grid} Bootstrap (\texttt{col-sm-*}, \texttt{col-lg-*})
para adaptar todas as telas ao viewport de 375 px (smartphone) sem overflow
horizontal.


% ################################################################
%  CAPÍTULO 3 (ou 4) — METODOLOGIA
% ################################################################

% ----------------------------------------------------------------
%  Subseção: Implementação da Interface — inserir em Metodologia
% ----------------------------------------------------------------
\subsection{Implementação da Camada de Interface}

O desenvolvimento da camada de interface seguiu um processo iterativo
composto por quatro etapas principais:

\begin{enumerate}
  \item \textbf{Diagnóstico} — inspeção de cada tela existente identificando
    deficiências visuais, HTML malformado e ausência de padrão entre
    componentes.
  \item \textbf{Priorização} — as correções foram ordenadas por impacto no
    usuário: primeiro os defeitos estruturais (HTML inválido, rotas
    quebradas), em seguida melhorias de usabilidade (formulários, tabelas,
    modais) e por último aprimoramentos visuais e adição de funcionalidades
    novas (gráfico de \textit{doughnut}).
  \item \textbf{Implementação} — as alterações foram realizadas arquivo a
    arquivo, com compilação Maven (\texttt{mvn compile}) após cada grupo de
    mudanças para garantir ausência de erros de \textit{template}.
  \item \textbf{Verificação} — cada tela foi verificada segundo o
    \textit{checklist} de testes de interface (Apêndice~A), que cobre os
    fluxos principais e os estados de erro em resolução \textit{desktop}
    (1280 px) e \textit{mobile} (375 px).
\end{enumerate}

\subsection{Controle de Qualidade Estático}

O projeto é analisado automaticamente pelo Qodana~\cite{qodana}, ferramenta
de análise estática da JetBrains integrada ao pipeline CI/CD do repositório
GitHub.  Durante o desenvolvimento foram identificados e corrigidos dois
grupos de avisos:

\begin{itemize}
  \item \textbf{Bulk operation} — substituição de laços \texttt{for} que
    populavam mapas linha a linha por chamadas \texttt{stream().collect
    (Collectors.toMap(...))} e \texttt{stream().collect
    (Collectors.groupingBy(...))}, conforme recomendação de estilo da
    plataforma Java.
  \item \textbf{Nullability / data flow} — substituição de comparador
    lambda com acesso encadeado sem verificação de nulo por
    \texttt{Comparator.comparing(AgendamentoResponse::getDataHora)},
    garantindo que a referência de método é nula-segura em tempo de compilação.
\end{itemize}

\subsection{Estratégia de Ramificação (\textit{Branching})}

O desenvolvimento seguiu o fluxo GitHub Flow: cada conjunto de alterações é
realizado em uma \textit{branch} de funcionalidade (\texttt{feature/}),
submetida como \textit{Pull Request} para revisão automatizada (Qodana) e
mesclada à \textit{branch} principal (\texttt{main}) após aprovação dos
\textit{checks} de CI.


% ################################################################
%  CAPÍTULO DE RESULTADOS E DISCUSSÃO
% ################################################################

\section{Resultados e Discussão}

Esta seção apresenta os resultados obtidos com a implementação da camada de
interface, analisando cada módulo do sistema sob os aspectos funcional,
visual e de usabilidade.

% ......................
\subsection{Painel Administrativo}

\subsubsection{Login}

A tela de login (\texttt{/web/login}, Figura~\ref{fig:tela-login}) passou a
exibir um fundo escuro com gradiente e um \textit{card} centralizado de alta
profundidade visual (\texttt{box-shadow: 0 24px 60px rgba(0,0,0,0.4)}).  O
campo de senha conta com um botão de alternância de visibilidade
(\textit{password toggle}), funcionalidade recomendada pelo NIST SP
800-63B~\cite{nist80063b} para redução de erros de autenticação.

\begin{figure}[H]
  \centering
  \includegraphics[width=0.55\textwidth]{figuras/tela-login.png}
  \caption{Tela de login com gradiente escuro e alternância de senha.}
  \label{fig:tela-login}
\end{figure}

\subsubsection{Navbar e Layout Base}

A barra de navegação (\texttt{layout/base.html}, Figura~\ref{fig:admin-navbar})
foi reconstruída com Bootstrap Navbar responsiva, ícones semânticos
Bootstrap Icons em cada item de menu, destaque dinâmico do item ativo por
JavaScript e colapso em \textit{hamburger menu} em telas menores que 992 px.
Alertas de sucesso e erro são exibidos com ícone e botão de fechamento
(\texttt{data-bs-dismiss="alert"}).

\begin{figure}[H]
  \centering
  \includegraphics[width=0.95\textwidth]{figuras/admin-navbar.png}
  \caption{Navbar administrativa com ícones e responsividade.}
  \label{fig:admin-navbar}
\end{figure}

\subsubsection{\textit{Dashboard}}

O \textit{dashboard} (\texttt{/dashboard}, Figura~\ref{fig:dashboard})
concentra as métricas operacionais mais relevantes em uma única tela:

\begin{itemize}
  \item \textbf{Quatro cartões de métricas} com fundo em gradiente, valor
    numérico em tipografia \textit{bold} de grande porte e efeito
    \textit{hover} de elevação (transform translateY(-4px)).
  \item \textbf{Gráfico de \textit{doughnut}} (Chart.js) com distribuição
    percentual dos agendamentos por status (Figura~\ref{fig:dashboard-chart}).
  \item \textbf{Barras de progresso Bootstrap} como visualização alternativa
    nativa do servidor (independente de JavaScript), exibindo contagem e
    percentual de cada status.
  \item \textbf{Tabela de próximos agendamentos} — até três registros com
    status AGENDADO ou CONFIRMADO, ordenados cronologicamente.
  \item \textbf{Acesso rápido} — três \textit{cards} com contadores de
    clientes, profissionais e serviços, com botões \textit{Gerenciar} e
    \textit{Novo} diretos.
\end{itemize}

\begin{figure}[H]
  \centering
  \includegraphics[width=0.95\textwidth]{figuras/dashboard.png}
  \caption{\textit{Dashboard} com métricas, gráfico e acesso rápido.}
  \label{fig:dashboard}
\end{figure}

\begin{figure}[H]
  \centering
  \includegraphics[width=0.55\textwidth]{figuras/dashboard-chart.png}
  \caption{Gráfico \textit{doughnut} de distribuição de agendamentos por status.}
  \label{fig:dashboard-chart}
\end{figure}

A Tabela~\ref{tab:cores-status} detalha a paleta de cores adotada para os
status, consistente entre o gráfico, as barras de progresso e os
\textit{badges} das tabelas.

\begin{table}[H]
\centering
\begin{tabular}{lll}
\hline
\textbf{Status} & \textbf{Cor}  & \textbf{Código Hexadecimal} \\ \hline
AGENDADO        & Azul-violeta  & \texttt{\#667eea}           \\
CONFIRMADO      & Verde         & \texttt{\#28a745}           \\
CONCLUÍDO       & Cinza         & \texttt{\#6c757d}           \\
CANCELADO       & Vermelho      & \texttt{\#dc3545}           \\ \hline
\end{tabular}
\caption{Paleta de cores dos status de agendamento.}
\label{tab:cores-status}
\end{table}

\paragraph{Desafio técnico — Chart.js em fragmento Thymeleaf.}
A integração do Chart.js com Thymeleaf exigiu a resolução de um problema
não documentado de forma proeminente: a diretiva \texttt{th:inline="javascript"}
não é processada quando o fragmento que a contém é passado como parâmetro a
um \textit{layout} template.  A solução implementada utilizou atributos
\texttt{data-*} no elemento \texttt{<canvas>} (injetados via \texttt{th:attr}
no servidor) e leitura via \texttt{canvas.dataset} no JavaScript do cliente,
conforme descrito na Listagem~\ref{lst:chartjs-dataset}.  Esta abordagem é
igualmente recomendada por questões de segurança: evita interpolação direta
de dados em contexto JavaScript, que poderia introduzir vulnerabilidades XSS.

\subsubsection{Listagens Administrativas}

Todas as telas de listagem seguem um padrão visual unificado
(Figuras~\ref{fig:clientes-lista},~\ref{fig:agendamentos-lista}):
cabeçalho com contador e botão de criação; busca em tempo real via JavaScript
(clientes, profissionais, serviços) ou filtro por formulário GET (agendamentos);
tabela Bootstrap \texttt{table-hover} com botões de ação ícone-apenas;
modal de confirmação para exclusão.

A listagem de agendamentos apresentava HTML malformado na versão anterior —
\texttt{<table>} sem classes Bootstrap, \texttt{</div>} órfãos e referências
a modais inexistentes.  A tela foi reescrita na íntegra, preservando toda a
lógica de transição de status e exclusão.

\begin{figure}[H]
  \centering
  \includegraphics[width=0.95\textwidth]{figuras/clientes-lista.png}
  \caption{Listagem de clientes com busca em tempo real e ações por ícone.}
  \label{fig:clientes-lista}
\end{figure}

\begin{figure}[H]
  \centering
  \includegraphics[width=0.95\textwidth]{figuras/agendamentos-lista.png}
  \caption{Listagem de agendamentos com filtros, badges de status e modais.}
  \label{fig:agendamentos-lista}
\end{figure}

\subsection{Portal do Cliente}

\subsubsection{\textit{Home} do Portal}

A página inicial do portal (\texttt{/portal}, Figura~\ref{fig:portal-home})
foi reformulada com um cartão \textit{hero} em gradiente
(\texttt{linear-gradient(135deg, \#667eea 0\%, \#764ba2 100\%)}), substituindo
o \textit{layout} simples anterior.  Três \textit{cards} de estatísticas
exibem serviços disponíveis, total de profissionais e agendamentos do cliente
autenticado.

\begin{figure}[H]
  \centering
  \includegraphics[width=0.95\textwidth]{figuras/portal-home.png}
  \caption{\textit{Home} do portal com \textit{hero card} e estatísticas.}
  \label{fig:portal-home}
\end{figure}

\subsubsection{\textit{Wizard} de Agendamento}

O fluxo de agendamento em quatro etapas
(Figuras~\ref{fig:wizard-step1},~\ref{fig:wizard-step3},~\ref{fig:wizard-step4})
guia o cliente desde a seleção do serviço até a confirmação com número de
protocolo.  O indicador de progresso visual (\textit{stepper}) comunica em
qual etapa o cliente se encontra e quais já foram concluídas (ícone de
verificação verde).  A Etapa 1 continha dois blocos \texttt{<script>}
duplicados que registravam os mesmos \textit{event listeners}; o bloco
redundante foi removido, eliminando a execução duplicada das funções de
seleção.

\begin{figure}[H]
  \centering
  \includegraphics[width=0.75\textwidth]{figuras/wizard-step1.png}
  \caption{Etapa 1 do \textit{wizard}: seleção de serviço por \textit{card}.}
  \label{fig:wizard-step1}
\end{figure}

\begin{figure}[H]
  \centering
  \includegraphics[width=0.65\textwidth]{figuras/wizard-step3.png}
  \caption{Etapa 3 do \textit{wizard}: revisão e confirmação do agendamento.}
  \label{fig:wizard-step3}
\end{figure}

\begin{figure}[H]
  \centering
  \includegraphics[width=0.65\textwidth]{figuras/wizard-step4.png}
  \caption{Etapa 4: tela de sucesso com protocolo de confirmação.}
  \label{fig:wizard-step4}
\end{figure}

\subsection{Melhorias de Usabilidade — Síntese}

A Tabela~\ref{tab:melhorias} resume as principais melhorias implementadas,
associando cada uma à heurística de Nielsen correspondente.

\begin{table}[H]
\centering
\small
\begin{tabular}{p{5cm} p{4cm} p{4.5cm}}
\hline
\textbf{Melhoria} & \textbf{Heurística} & \textbf{Resultado} \\ \hline
Fundo gradiente + card sombra na tela de login &
  Aesthetic \& minimalist design &
  Foco no formulário, redução de distração \\[4pt]
Ícones nos itens de menu &
  Recognition rather than recall &
  Navegação mais rápida \\[4pt]
\textit{Password toggle} &
  Error prevention &
  Menos erros de digitação de senha \\[4pt]
Modal de confirmação para exclusão &
  User control and freedom &
  Reversibilidade de ações destrutivas \\[4pt]
Gráfico + barras de progresso no dashboard &
  Visibility of system status &
  Distribuição de agendamentos em tempo real \\[4pt]
Busca em tempo real nas listagens &
  Efficiency of use &
  Localização rápida de registros \\[4pt]
Correção HTML malformado (agendamentos) &
  Consistency and standards &
  Comportamento previsível dos componentes \\ \hline
\end{tabular}
\caption{Melhorias de usabilidade e heurísticas de Nielsen correspondentes.}
\label{tab:melhorias}
\end{table}

\subsection{Discussão}

Os resultados demonstram que a adoção de um \textit{design system} baseado
em variáveis CSS globais e nos componentes Bootstrap 5 foi fundamental para
manter a consistência visual entre as mais de quinze telas do sistema, sem
impacto na lógica de negócio.  A separação entre camada de apresentação
(Thymeleaf + Bootstrap) e camada de negócio (serviços e repositórios Spring)
permaneceu preservada em todas as alterações realizadas.

A decisão de oferecer barras de progresso Bootstrap como alternativa nativa
ao gráfico Chart.js mostrou-se acertada sob dois aspectos: primeiro, garante
que o \textit{dashboard} exiba informações de status mesmo em ambientes com
JavaScript desabilitado ou bloqueado; segundo, evidencia os dados numéricos
de forma inequívoca, complementando a visualização proporcional do gráfico.

O desafio técnico da integração Chart.js/Thymeleaf ilustra um aspecto
relevante do desenvolvimento com \textit{server-side rendering}: a
necessidade de compreender os limites do processamento de \textit{templates}
no lado do servidor antes de tomar decisões de integração com bibliotecas
JavaScript.  A solução por atributos \texttt{data-*} é idiomática na web
moderna e agnóstica ao \textit{framework} de \textit{templates}, podendo ser
reaproveitada em projetos futuros que utilizem Freemarker, Pebble ou JSP.


% ################################################################
%  CAPÍTULO — CONCLUSÃO
% ################################################################

\section{Conclusão}

Este trabalho apresentou o desenvolvimento de um sistema de agendamento
construído sobre a plataforma Spring Boot, com ênfase na qualidade da camada
de interface e na experiência do usuário.  O sistema atende ao objetivo
principal de oferecer um painel administrativo completo e um portal do
cliente com fluxo de agendamento guiado, ambos responsivos e visualmente
adequados à apresentação acadêmica e ao ambiente de produção.

\subsection{Contribuições}

As principais contribuições técnicas do trabalho são:

\begin{itemize}
  \item \textbf{Arquitetura de \textit{templates} Thymeleaf} com fragmentos e
    layouts-base reutilizáveis, demonstrando boas práticas de organização de
    \textit{templates} em projetos Spring Boot.
  \item \textbf{Padrão de integração Chart.js/Thymeleaf} via atributos
    \texttt{data-*}, documentando e resolvendo a incompatibilidade entre
    \texttt{th:inline="javascript"} e fragmentos parametrizados, que não
    consta de forma destacada na documentação oficial do Thymeleaf.
  \item \textbf{\textit{Design system} coerente} — paleta, espaçamentos, raios
    de borda e sombras definidos em variáveis CSS (\texttt{:root}), facilmente
    adaptável a outras identidades visuais.
  \item \textbf{Fluxo de agendamento em \textit{wizard}} com indicador de
    progresso, validação por etapa e tela de confirmação com protocolo,
    padrão aplicável a outros domínios de negócio.
  \item \textbf{Pipeline de qualidade estática} — correção de avisos Qodana
    sobre operações de coleção e nulabilidade, demonstrando integração de
    análise estática no fluxo de desenvolvimento.
\end{itemize}

\subsection{Limitações}

O trabalho apresenta as seguintes limitações que podem ser endereçadas em
trabalhos futuros:

\begin{itemize}
  \item \textbf{Testes de interface automatizados} — o \textit{checklist} de
    testes foi realizado manualmente.  A automação com Selenium ou Playwright
    aumentaria a cobertura e a rastreabilidade das verificações.
  \item \textbf{Cobertura de testes de unidade} — os testes automatizados
    existentes cobrem os serviços principais; a camada de \textit{controllers}
    não possui testes de integração com MockMvc.
  \item \textbf{Acessibilidade formal} — embora as interfaces sigam
    recomendações gerais de contraste WCAG 2.1, não foi realizada auditoria
    formal com ferramentas como Axe ou Lighthouse.
  \item \textbf{Internacionalização} — o sistema está fixado em português
    brasileiro; suporte a múltiplos idiomas exigiria a adoção do mecanismo
    \texttt{MessageSource} do Spring.
\end{itemize}

\subsection{Trabalhos Futuros}

Como extensões naturais do sistema, destacam-se:

\begin{itemize}
  \item Implementação de notificações por e-mail e/ou WhatsApp para lembrete
    de agendamentos com proximidade da data.
  \item Exposição de uma API REST com autenticação JWT para integração com
    aplicativos móveis nativos.
  \item Adição de relatórios exportáveis (PDF/CSV) de agendamentos por
    período, profissional ou serviço.
  \item Suporte a múltiplos estabelecimentos (multi-tenant) na mesma
    instalação.
\end{itemize}

Em síntese, o sistema de agendamento desenvolvido neste trabalho cumpre seu
objetivo acadêmico e funcional, entregando uma plataforma operacional com
interfaces de qualidade profissional, arquitetura limpa e pipeline de
qualidade integrado.


% ################################################################
%  REFERÊNCIAS — adicionar ao arquivo .bib do projeto
% ################################################################

%
% @misc{springboot,
%   author = {{Pivotal Software / VMware}},
%   title  = {Spring Boot Reference Documentation},
%   year   = {2024},
%   url    = {https://docs.spring.io/spring-boot/docs/current/reference/html}
% }
% @misc{springsecurity,
%   author = {{Spring Security Team}},
%   title  = {Spring Security Reference},
%   year   = {2024},
%   url    = {https://docs.spring.io/spring-security/reference}
% }
% @misc{springdatajpa,
%   author = {{Spring Data Team}},
%   title  = {Spring Data JPA Reference},
%   year   = {2024},
%   url    = {https://docs.spring.io/spring-data/jpa/docs/current/reference/html}
% }
% @misc{thymeleaf,
%   author = {{Thymeleaf Team}},
%   title  = {Thymeleaf 3.x Documentation},
%   year   = {2024},
%   url    = {https://www.thymeleaf.org/documentation.html}
% }
% @misc{bootstrap5,
%   author = {{The Bootstrap Authors}},
%   title  = {Bootstrap 5.3 Documentation},
%   year   = {2024},
%   url    = {https://getbootstrap.com/docs/5.3}
% }
% @misc{bootstrapicons,
%   author = {{The Bootstrap Authors}},
%   title  = {Bootstrap Icons 1.11},
%   year   = {2024},
%   url    = {https://icons.getbootstrap.com}
% }
% @misc{chartjs,
%   author = {{Chart.js Contributors}},
%   title  = {Chart.js 4.x Documentation},
%   year   = {2024},
%   url    = {https://www.chartjs.org/docs/latest}
% }
% @misc{qodana,
%   author = {{JetBrains}},
%   title  = {Qodana — Static Code Analysis},
%   year   = {2024},
%   url    = {https://www.jetbrains.com/qodana}
% }
% @book{nielsen1994,
%   author    = {Nielsen, Jakob},
%   title     = {Usability Engineering},
%   publisher = {Morgan Kaufmann},
%   year      = {1994}
% }
% @techreport{wcag21,
%   author      = {{W3C Web Content Accessibility Guidelines Working Group}},
%   title       = {{WCAG} 2.1},
%   institution = {W3C},
%   year        = {2018},
%   url         = {https://www.w3.org/TR/WCAG21}
% }
% @techreport{nist80063b,
%   author      = {Grassi, Paul A. and others},
%   title       = {{NIST SP 800-63B}: Digital Identity Guidelines -- Authentication},
%   institution = {NIST},
%   year        = {2017}
% }
% @article{pierce2004,
%   author  = {Pierce, Jon L.},
%   title   = {The Role of Icons in User Interface Design},
%   journal = {Journal of Usability Studies},
%   year    = {2004}
% }
% @book{gamma1994,
%   author    = {Gamma, Erich and Helm, Richard and Johnson, Ralph and Vlissides, John},
%   title     = {Design Patterns: Elements of Reusable Object-Oriented Software},
%   publisher = {Addison-Wesley},
%   year      = {1994}
% }
.
% ================================================================


% ################################################################
%  CAPÍTULO 1 — INTRODUÇÃO
%  (insira na seção de motivação / objetivos do capítulo 1)
% ################################################################

% ----------------------------------------------------------------
%  Parágrafo de contextualização — inserir após a apresentação do
%  problema de pesquisa
% ----------------------------------------------------------------
A interface de uma aplicação web é o primeiro ponto de contato entre
o usuário e o sistema; uma interface mal projetada pode comprometer a
adoção da ferramenta independentemente da qualidade do seu código de
negócio~\cite{nielsen1994}.  Nesse contexto, o presente trabalho dedica
atenção especial ao desenvolvimento de uma camada de apresentação
profissional, responsiva e acessível, utilizando tecnologias consolidadas
do ecossistema Java — Thymeleaf, Bootstrap 5 e Chart.js — sobre a
plataforma Spring Boot.

% ----------------------------------------------------------------
%  Objetivo específico — adicionar à lista de objetivos específicos
% ----------------------------------------------------------------
\begin{itemize}
  \item Projetar e implementar interfaces administrativas e um portal do
    cliente com padrões visuais modernos, responsividade \textit{mobile-first}
    e navegação intuitiva, adequadas à apresentação acadêmica e ao uso
    em ambiente de produção.
  \item Incorporar visualização de dados (\textit{dashboard} com gráfico de
    \textit{doughnut} e barras de progresso) que permita ao gestor monitorar
    a distribuição de agendamentos por status em tempo real.
\end{itemize}

% ----------------------------------------------------------------
%  Justificativa — inserir no parágrafo de justificativa
% ----------------------------------------------------------------
A escolha por \textit{server-side rendering} com Thymeleaf, em vez de um
\textit{framework} SPA (\textit{Single Page Application}), justifica-se pela
natureza do projeto: um sistema de agendamento com domínio bem definido,
cuja maioria das interações do usuário segue fluxos lineares (cadastro,
listagem, agendamento). Essa abordagem reduz a complexidade de implantação,
elimina a necessidade de uma API REST separada para as telas internas e
facilita a validação de formulários no servidor com o Spring MVC
\textit{BindingResult}.


% ################################################################
%  CAPÍTULO 2 — FUNDAMENTAÇÃO TEÓRICA
% ################################################################

% ----------------------------------------------------------------
%  Seção: Tecnologias de Backend
%  (inserir na subseção de Spring Boot / Spring MVC)
% ----------------------------------------------------------------
\subsection{Spring Boot e Spring MVC}

O Spring Boot~\cite{springboot} é um \textit{framework} que simplifica a
configuração e a implantação de aplicações Java baseadas no ecossistema
Spring.  Ele adota o paradigma de \textit{convention over configuration}:
por meio de dependências \textit{starter} e auto-configurações, o
desenvolvedor pode iniciar uma aplicação web produtiva com poucas linhas de
configuração.

O Spring MVC implementa o padrão arquitetural
\textit{Model-View-Controller}~\cite{gamma1994}: o \texttt{@Controller}
recebe a requisição HTTP, delega a lógica de negócio à camada de serviço e
devolve um nome de \textit{view} ao \texttt{DispatcherServlet}, que aciona
o motor de \textit{templates} Thymeleaf para gerar o HTML final enviado ao
navegador.

\subsection{Spring Security}

O Spring Security~\cite{springsecurity} fornece autenticação e autorização ao
sistema.  No contexto do painel administrativo, a autenticação é realizada
por meio de formulário HTML com sessão gerenciada pelo servidor; o
\texttt{SecurityConfig} configura as rotas protegidas (\texttt{/dashboard},
\texttt{/web/**}) e as rotas públicas (\texttt{/portal/**}).  O portal do
cliente utiliza autenticação manual via \texttt{HttpSession}, com senhas
armazenadas em \textit{hash} BCrypt (\texttt{BCryptPasswordEncoder}).

\subsection{Spring Data JPA e PostgreSQL}

A persistência é gerenciada pelo Spring Data JPA~\cite{springdatajpa},
que abstrai as operações de banco de dados por meio de interfaces
\texttt{JpaRepository}.  As entidades — \texttt{Cliente},
\texttt{Profissional}, \texttt{Servico} e \texttt{Agendamento} — são
mapeadas com anotações JPA (\texttt{@Entity}, \texttt{@ManyToOne},
\texttt{@Enumerated}) e o banco de dados utilizado é o PostgreSQL.

% ----------------------------------------------------------------
%  Seção: Tecnologias de Interface e Frontend
%  (inserir como subseção de Tecnologias Utilizadas)
% ----------------------------------------------------------------
\subsection{Tecnologias de Interface e \textit{Frontend}}

A camada de apresentação foi construída inteiramente sobre tecnologias
consolidadas do ecossistema Java, garantindo compatibilidade nativa com o
Spring Boot, manutenibilidade e facilidade de evolução futura.

\subsubsection{Thymeleaf}

O Thymeleaf~\cite{thymeleaf} é um motor de \textit{templates} para a
plataforma Java integrado nativamente ao Spring MVC.  Toda a geração de
HTML é realizada no servidor (\textit{server-side rendering}), eliminando a
necessidade de APIs REST separadas para as telas administrativas e o portal
do cliente.  Os \textit{templates} herdam de layouts-base por meio dos
mecanismos de fragmento (\texttt{th:replace} / \texttt{th:fragment}),
garantindo consistência visual e evitando duplicação de código.

Um aspecto técnico relevante explorado neste trabalho diz respeito à
interação entre fragmentos Thymeleaf e JavaScript: quando um fragmento é
passado como parâmetro a um \textit{layout} via
\texttt{th:replace="\textasciitilde\{layout :: html(\textasciitilde\{:: conteudo\})\}"},
qualquer diretiva \texttt{th:inline="javascript"} dentro desse fragmento
não é processada de forma confiável~\cite{thymeleaf}.  A solução adotada
foi injetar os dados do controlador como atributos \texttt{data-*} no
elemento HTML (\texttt{th:attr}) e lê-los no JavaScript puro por meio de
\texttt{element.dataset}, eliminando a dependência do processamento
\textit{inline} e mitigando riscos de XSS por interpolação direta.

\subsubsection{Bootstrap 5}

O Bootstrap 5.3.3~\cite{bootstrap5} fornece o sistema de \textit{grid}
responsivo baseado em \textit{Flexbox}, componentes visuais (Navbar, Modal,
Card, Badge, Table, Progress) e utilitários CSS utilizados em todas as
telas.  A versão 5 eliminou a dependência do jQuery, tornando o
\textit{framework} mais leve e compatível com o JavaScript puro utilizado
nos \textit{scripts} de validação e interação.

\subsubsection{Bootstrap Icons}

A biblioteca Bootstrap Icons 1.11.3~\cite{bootstrapicons} é empregada como
sistema de ícones SVG (\textit{icon font}), presente em botões, títulos de
página, itens de menu e \textit{badges} de status.  O uso consistente de
ícones contextuais acelera o reconhecimento de destinos de navegação e
reduz a carga cognitiva do usuário~\cite{pierce2004}.

\subsubsection{Chart.js}

O Chart.js 4.4.3~\cite{chartjs} é incorporado via CDN para renderizar o
gráfico de \textit{doughnut} de distribuição de agendamentos por status no
\textit{dashboard} administrativo.  Os dados são injetados pelo Thymeleaf
por meio de atributos \texttt{data-*} no elemento \texttt{<canvas>}
(Listagem~\ref{lst:chartjs-dataset}) e lidos no JavaScript por meio de
\texttt{canvas.dataset}, garantindo separação clara entre lógica de
servidor e apresentação no cliente.

\begin{lstlisting}[language=HTML, caption={Injeção de dados via atributos \texttt{data-*}
  e inicialização do Chart.js via \texttt{canvas.dataset}.},
  label=lst:chartjs-dataset, basicstyle=\small\ttfamily]
<!-- Template Thymeleaf: injeta dados no canvas -->
<canvas id="graficoAgendamentosStatus"
        th:attr="data-agendado=${totalAgendado},
                 data-confirmado=${totalConfirmado},
                 data-concluido=${totalConcluido},
                 data-cancelado=${totalCancelado}"
        style="max-height:220px;"></canvas>

<!-- JavaScript puro: le os dados sem th:inline -->
<script>
(function () {
    var c = document.getElementById('graficoAgendamentosStatus');
    var agendado   = parseInt(c.dataset.agendado   || '0', 10);
    var confirmado = parseInt(c.dataset.confirmado || '0', 10);
    var concluido  = parseInt(c.dataset.concluido  || '0', 10);
    var cancelado  = parseInt(c.dataset.cancelado  || '0', 10);
    new Chart(c, { type: 'doughnut',
        data: { labels: ['AGENDADO','CONFIRMADO','CONCLUIDO','CANCELADO'],
            datasets: [{ data: [agendado, confirmado, concluido, cancelado],
                backgroundColor: ['#667eea','#28a745','#6c757d','#dc3545'],
                borderWidth: 2, borderColor: '#fff' }] },
        options: { plugins: { legend: { display: false } },
            cutout: '65%', responsive: true,
            maintainAspectRatio: false } });
})();
</script>
\end{lstlisting}

% ----------------------------------------------------------------
%  Seção: Usabilidade e Design de Interface
% ----------------------------------------------------------------
\subsection{Usabilidade e \textit{Design} de Interface}

\subsubsection{Heurísticas de Nielsen}

Nielsen~\cite{nielsen1994} definiu dez heurísticas de usabilidade para o
\textit{design} de interfaces de usuário, amplamente adotadas como
referência na avaliação e no projeto de sistemas interativos.  Dentre as
heurísticas mais diretamente aplicadas neste trabalho destacam-se:
\textit{Consistency and Standards} (uso uniforme de cores, ícones e padrões
de interação), \textit{Recognition rather than recall} (ícones contextuais
nos menus) e \textit{User control and freedom} (botões de cancelar/voltar
em todos os formulários e modais de confirmação para ações destrutivas).

\subsubsection{Responsividade \textit{Mobile-First}}

A abordagem \textit{mobile-first} preconiza que o \textit{layout} base seja
projetado para telas pequenas e progressivamente aprimorado para telas
maiores por meio de \textit{media queries}~\cite{bootstrap5}.  O sistema
emprega o \textit{grid} Bootstrap (\texttt{col-sm-*}, \texttt{col-lg-*})
para adaptar todas as telas ao viewport de 375 px (smartphone) sem overflow
horizontal.


% ################################################################
%  CAPÍTULO 3 (ou 4) — METODOLOGIA
% ################################################################

% ----------------------------------------------------------------
%  Subseção: Implementação da Interface — inserir em Metodologia
% ----------------------------------------------------------------
\subsection{Implementação da Camada de Interface}

O desenvolvimento da camada de interface seguiu um processo iterativo
composto por quatro etapas principais:

\begin{enumerate}
  \item \textbf{Diagnóstico} — inspeção de cada tela existente identificando
    deficiências visuais, HTML malformado e ausência de padrão entre
    componentes.
  \item \textbf{Priorização} — as correções foram ordenadas por impacto no
    usuário: primeiro os defeitos estruturais (HTML inválido, rotas
    quebradas), em seguida melhorias de usabilidade (formulários, tabelas,
    modais) e por último aprimoramentos visuais e adição de funcionalidades
    novas (gráfico de \textit{doughnut}).
  \item \textbf{Implementação} — as alterações foram realizadas arquivo a
    arquivo, com compilação Maven (\texttt{mvn compile}) após cada grupo de
    mudanças para garantir ausência de erros de \textit{template}.
  \item \textbf{Verificação} — cada tela foi verificada segundo o
    \textit{checklist} de testes de interface (Apêndice~A), que cobre os
    fluxos principais e os estados de erro em resolução \textit{desktop}
    (1280 px) e \textit{mobile} (375 px).
\end{enumerate}

\subsection{Controle de Qualidade Estático}

O projeto é analisado automaticamente pelo Qodana~\cite{qodana}, ferramenta
de análise estática da JetBrains integrada ao pipeline CI/CD do repositório
GitHub.  Durante o desenvolvimento foram identificados e corrigidos dois
grupos de avisos:

\begin{itemize}
  \item \textbf{Bulk operation} — substituição de laços \texttt{for} que
    populavam mapas linha a linha por chamadas \texttt{stream().collect
    (Collectors.toMap(...))} e \texttt{stream().collect
    (Collectors.groupingBy(...))}, conforme recomendação de estilo da
    plataforma Java.
  \item \textbf{Nullability / data flow} — substituição de comparador
    lambda com acesso encadeado sem verificação de nulo por
    \texttt{Comparator.comparing(AgendamentoResponse::getDataHora)},
    garantindo que a referência de método é nula-segura em tempo de compilação.
\end{itemize}

\subsection{Estratégia de Ramificação (\textit{Branching})}

O desenvolvimento seguiu o fluxo GitHub Flow: cada conjunto de alterações é
realizado em uma \textit{branch} de funcionalidade (\texttt{feature/}),
submetida como \textit{Pull Request} para revisão automatizada (Qodana) e
mesclada à \textit{branch} principal (\texttt{main}) após aprovação dos
\textit{checks} de CI.


% ################################################################
%  CAPÍTULO DE RESULTADOS E DISCUSSÃO
% ################################################################

\section{Resultados e Discussão}

Esta seção apresenta os resultados obtidos com a implementação da camada de
interface, analisando cada módulo do sistema sob os aspectos funcional,
visual e de usabilidade.

% ......................
\subsection{Painel Administrativo}

\subsubsection{Login}

A tela de login (\texttt{/web/login}, Figura~\ref{fig:tela-login}) passou a
exibir um fundo escuro com gradiente e um \textit{card} centralizado de alta
profundidade visual (\texttt{box-shadow: 0 24px 60px rgba(0,0,0,0.4)}).  O
campo de senha conta com um botão de alternância de visibilidade
(\textit{password toggle}), funcionalidade recomendada pelo NIST SP
800-63B~\cite{nist80063b} para redução de erros de autenticação.

\begin{figure}[H]
  \centering
  \includegraphics[width=0.55\textwidth]{figuras/tela-login.png}
  \caption{Tela de login com gradiente escuro e alternância de senha.}
  \label{fig:tela-login}
\end{figure}

\subsubsection{Navbar e Layout Base}

A barra de navegação (\texttt{layout/base.html}, Figura~\ref{fig:admin-navbar})
foi reconstruída com Bootstrap Navbar responsiva, ícones semânticos
Bootstrap Icons em cada item de menu, destaque dinâmico do item ativo por
JavaScript e colapso em \textit{hamburger menu} em telas menores que 992 px.
Alertas de sucesso e erro são exibidos com ícone e botão de fechamento
(\texttt{data-bs-dismiss="alert"}).

\begin{figure}[H]
  \centering
  \includegraphics[width=0.95\textwidth]{figuras/admin-navbar.png}
  \caption{Navbar administrativa com ícones e responsividade.}
  \label{fig:admin-navbar}
\end{figure}

\subsubsection{\textit{Dashboard}}

O \textit{dashboard} (\texttt{/dashboard}, Figura~\ref{fig:dashboard})
concentra as métricas operacionais mais relevantes em uma única tela:

\begin{itemize}
  \item \textbf{Quatro cartões de métricas} com fundo em gradiente, valor
    numérico em tipografia \textit{bold} de grande porte e efeito
    \textit{hover} de elevação (transform translateY(-4px)).
  \item \textbf{Gráfico de \textit{doughnut}} (Chart.js) com distribuição
    percentual dos agendamentos por status (Figura~\ref{fig:dashboard-chart}).
  \item \textbf{Barras de progresso Bootstrap} como visualização alternativa
    nativa do servidor (independente de JavaScript), exibindo contagem e
    percentual de cada status.
  \item \textbf{Tabela de próximos agendamentos} — até três registros com
    status AGENDADO ou CONFIRMADO, ordenados cronologicamente.
  \item \textbf{Acesso rápido} — três \textit{cards} com contadores de
    clientes, profissionais e serviços, com botões \textit{Gerenciar} e
    \textit{Novo} diretos.
\end{itemize}

\begin{figure}[H]
  \centering
  \includegraphics[width=0.95\textwidth]{figuras/dashboard.png}
  \caption{\textit{Dashboard} com métricas, gráfico e acesso rápido.}
  \label{fig:dashboard}
\end{figure}

\begin{figure}[H]
  \centering
  \includegraphics[width=0.55\textwidth]{figuras/dashboard-chart.png}
  \caption{Gráfico \textit{doughnut} de distribuição de agendamentos por status.}
  \label{fig:dashboard-chart}
\end{figure}

A Tabela~\ref{tab:cores-status} detalha a paleta de cores adotada para os
status, consistente entre o gráfico, as barras de progresso e os
\textit{badges} das tabelas.

\begin{table}[H]
\centering
\begin{tabular}{lll}
\hline
\textbf{Status} & \textbf{Cor}  & \textbf{Código Hexadecimal} \\ \hline
AGENDADO        & Azul-violeta  & \texttt{\#667eea}           \\
CONFIRMADO      & Verde         & \texttt{\#28a745}           \\
CONCLUÍDO       & Cinza         & \texttt{\#6c757d}           \\
CANCELADO       & Vermelho      & \texttt{\#dc3545}           \\ \hline
\end{tabular}
\caption{Paleta de cores dos status de agendamento.}
\label{tab:cores-status}
\end{table}

\paragraph{Desafio técnico — Chart.js em fragmento Thymeleaf.}
A integração do Chart.js com Thymeleaf exigiu a resolução de um problema
não documentado de forma proeminente: a diretiva \texttt{th:inline="javascript"}
não é processada quando o fragmento que a contém é passado como parâmetro a
um \textit{layout} template.  A solução implementada utilizou atributos
\texttt{data-*} no elemento \texttt{<canvas>} (injetados via \texttt{th:attr}
no servidor) e leitura via \texttt{canvas.dataset} no JavaScript do cliente,
conforme descrito na Listagem~\ref{lst:chartjs-dataset}.  Esta abordagem é
igualmente recomendada por questões de segurança: evita interpolação direta
de dados em contexto JavaScript, que poderia introduzir vulnerabilidades XSS.

\subsubsection{Listagens Administrativas}

Todas as telas de listagem seguem um padrão visual unificado
(Figuras~\ref{fig:clientes-lista},~\ref{fig:agendamentos-lista}):
cabeçalho com contador e botão de criação; busca em tempo real via JavaScript
(clientes, profissionais, serviços) ou filtro por formulário GET (agendamentos);
tabela Bootstrap \texttt{table-hover} com botões de ação ícone-apenas;
modal de confirmação para exclusão.

A listagem de agendamentos apresentava HTML malformado na versão anterior —
\texttt{<table>} sem classes Bootstrap, \texttt{</div>} órfãos e referências
a modais inexistentes.  A tela foi reescrita na íntegra, preservando toda a
lógica de transição de status e exclusão.

\begin{figure}[H]
  \centering
  \includegraphics[width=0.95\textwidth]{figuras/clientes-lista.png}
  \caption{Listagem de clientes com busca em tempo real e ações por ícone.}
  \label{fig:clientes-lista}
\end{figure}

\begin{figure}[H]
  \centering
  \includegraphics[width=0.95\textwidth]{figuras/agendamentos-lista.png}
  \caption{Listagem de agendamentos com filtros, badges de status e modais.}
  \label{fig:agendamentos-lista}
\end{figure}

\subsection{Portal do Cliente}

\subsubsection{\textit{Home} do Portal}

A página inicial do portal (\texttt{/portal}, Figura~\ref{fig:portal-home})
foi reformulada com um cartão \textit{hero} em gradiente
(\texttt{linear-gradient(135deg, \#667eea 0\%, \#764ba2 100\%)}), substituindo
o \textit{layout} simples anterior.  Três \textit{cards} de estatísticas
exibem serviços disponíveis, total de profissionais e agendamentos do cliente
autenticado.

\begin{figure}[H]
  \centering
  \includegraphics[width=0.95\textwidth]{figuras/portal-home.png}
  \caption{\textit{Home} do portal com \textit{hero card} e estatísticas.}
  \label{fig:portal-home}
\end{figure}

\subsubsection{\textit{Wizard} de Agendamento}

O fluxo de agendamento em quatro etapas
(Figuras~\ref{fig:wizard-step1},~\ref{fig:wizard-step3},~\ref{fig:wizard-step4})
guia o cliente desde a seleção do serviço até a confirmação com número de
protocolo.  O indicador de progresso visual (\textit{stepper}) comunica em
qual etapa o cliente se encontra e quais já foram concluídas (ícone de
verificação verde).  A Etapa 1 continha dois blocos \texttt{<script>}
duplicados que registravam os mesmos \textit{event listeners}; o bloco
redundante foi removido, eliminando a execução duplicada das funções de
seleção.

\begin{figure}[H]
  \centering
  \includegraphics[width=0.75\textwidth]{figuras/wizard-step1.png}
  \caption{Etapa 1 do \textit{wizard}: seleção de serviço por \textit{card}.}
  \label{fig:wizard-step1}
\end{figure}

\begin{figure}[H]
  \centering
  \includegraphics[width=0.65\textwidth]{figuras/wizard-step3.png}
  \caption{Etapa 3 do \textit{wizard}: revisão e confirmação do agendamento.}
  \label{fig:wizard-step3}
\end{figure}

\begin{figure}[H]
  \centering
  \includegraphics[width=0.65\textwidth]{figuras/wizard-step4.png}
  \caption{Etapa 4: tela de sucesso com protocolo de confirmação.}
  \label{fig:wizard-step4}
\end{figure}

\subsection{Melhorias de Usabilidade — Síntese}

A Tabela~\ref{tab:melhorias} resume as principais melhorias implementadas,
associando cada uma à heurística de Nielsen correspondente.

\begin{table}[H]
\centering
\small
\begin{tabular}{p{5cm} p{4cm} p{4.5cm}}
\hline
\textbf{Melhoria} & \textbf{Heurística} & \textbf{Resultado} \\ \hline
Fundo gradiente + card sombra na tela de login &
  Aesthetic \& minimalist design &
  Foco no formulário, redução de distração \\[4pt]
Ícones nos itens de menu &
  Recognition rather than recall &
  Navegação mais rápida \\[4pt]
\textit{Password toggle} &
  Error prevention &
  Menos erros de digitação de senha \\[4pt]
Modal de confirmação para exclusão &
  User control and freedom &
  Reversibilidade de ações destrutivas \\[4pt]
Gráfico + barras de progresso no dashboard &
  Visibility of system status &
  Distribuição de agendamentos em tempo real \\[4pt]
Busca em tempo real nas listagens &
  Efficiency of use &
  Localização rápida de registros \\[4pt]
Correção HTML malformado (agendamentos) &
  Consistency and standards &
  Comportamento previsível dos componentes \\ \hline
\end{tabular}
\caption{Melhorias de usabilidade e heurísticas de Nielsen correspondentes.}
\label{tab:melhorias}
\end{table}

\subsection{Discussão}

Os resultados demonstram que a adoção de um \textit{design system} baseado
em variáveis CSS globais e nos componentes Bootstrap 5 foi fundamental para
manter a consistência visual entre as mais de quinze telas do sistema, sem
impacto na lógica de negócio.  A separação entre camada de apresentação
(Thymeleaf + Bootstrap) e camada de negócio (serviços e repositórios Spring)
permaneceu preservada em todas as alterações realizadas.

A decisão de oferecer barras de progresso Bootstrap como alternativa nativa
ao gráfico Chart.js mostrou-se acertada sob dois aspectos: primeiro, garante
que o \textit{dashboard} exiba informações de status mesmo em ambientes com
JavaScript desabilitado ou bloqueado; segundo, evidencia os dados numéricos
de forma inequívoca, complementando a visualização proporcional do gráfico.

O desafio técnico da integração Chart.js/Thymeleaf ilustra um aspecto
relevante do desenvolvimento com \textit{server-side rendering}: a
necessidade de compreender os limites do processamento de \textit{templates}
no lado do servidor antes de tomar decisões de integração com bibliotecas
JavaScript.  A solução por atributos \texttt{data-*} é idiomática na web
moderna e agnóstica ao \textit{framework} de \textit{templates}, podendo ser
reaproveitada em projetos futuros que utilizem Freemarker, Pebble ou JSP.


% ################################################################
%  CAPÍTULO — CONCLUSÃO
% ################################################################

\section{Conclusão}

Este trabalho apresentou o desenvolvimento de um sistema de agendamento
construído sobre a plataforma Spring Boot, com ênfase na qualidade da camada
de interface e na experiência do usuário.  O sistema atende ao objetivo
principal de oferecer um painel administrativo completo e um portal do
cliente com fluxo de agendamento guiado, ambos responsivos e visualmente
adequados à apresentação acadêmica e ao ambiente de produção.

\subsection{Contribuições}

As principais contribuições técnicas do trabalho são:

\begin{itemize}
  \item \textbf{Arquitetura de \textit{templates} Thymeleaf} com fragmentos e
    layouts-base reutilizáveis, demonstrando boas práticas de organização de
    \textit{templates} em projetos Spring Boot.
  \item \textbf{Padrão de integração Chart.js/Thymeleaf} via atributos
    \texttt{data-*}, documentando e resolvendo a incompatibilidade entre
    \texttt{th:inline="javascript"} e fragmentos parametrizados, que não
    consta de forma destacada na documentação oficial do Thymeleaf.
  \item \textbf{\textit{Design system} coerente} — paleta, espaçamentos, raios
    de borda e sombras definidos em variáveis CSS (\texttt{:root}), facilmente
    adaptável a outras identidades visuais.
  \item \textbf{Fluxo de agendamento em \textit{wizard}} com indicador de
    progresso, validação por etapa e tela de confirmação com protocolo,
    padrão aplicável a outros domínios de negócio.
  \item \textbf{Pipeline de qualidade estática} — correção de avisos Qodana
    sobre operações de coleção e nulabilidade, demonstrando integração de
    análise estática no fluxo de desenvolvimento.
\end{itemize}

\subsection{Limitações}

O trabalho apresenta as seguintes limitações que podem ser endereçadas em
trabalhos futuros:

\begin{itemize}
  \item \textbf{Testes de interface automatizados} — o \textit{checklist} de
    testes foi realizado manualmente.  A automação com Selenium ou Playwright
    aumentaria a cobertura e a rastreabilidade das verificações.
  \item \textbf{Cobertura de testes de unidade} — os testes automatizados
    existentes cobrem os serviços principais; a camada de \textit{controllers}
    não possui testes de integração com MockMvc.
  \item \textbf{Acessibilidade formal} — embora as interfaces sigam
    recomendações gerais de contraste WCAG 2.1, não foi realizada auditoria
    formal com ferramentas como Axe ou Lighthouse.
  \item \textbf{Internacionalização} — o sistema está fixado em português
    brasileiro; suporte a múltiplos idiomas exigiria a adoção do mecanismo
    \texttt{MessageSource} do Spring.
\end{itemize}

\subsection{Trabalhos Futuros}

Como extensões naturais do sistema, destacam-se:

\begin{itemize}
  \item Implementação de notificações por e-mail e/ou WhatsApp para lembrete
    de agendamentos com proximidade da data.
  \item Exposição de uma API REST com autenticação JWT para integração com
    aplicativos móveis nativos.
  \item Adição de relatórios exportáveis (PDF/CSV) de agendamentos por
    período, profissional ou serviço.
  \item Suporte a múltiplos estabelecimentos (multi-tenant) na mesma
    instalação.
\end{itemize}

Em síntese, o sistema de agendamento desenvolvido neste trabalho cumpre seu
objetivo acadêmico e funcional, entregando uma plataforma operacional com
interfaces de qualidade profissional, arquitetura limpa e pipeline de
qualidade integrado.


% ################################################################
%  REFERÊNCIAS — adicionar ao arquivo .bib do projeto
% ################################################################

%
% @misc{springboot,
%   author = {{Pivotal Software / VMware}},
%   title  = {Spring Boot Reference Documentation},
%   year   = {2024},
%   url    = {https://docs.spring.io/spring-boot/docs/current/reference/html}
% }
% @misc{springsecurity,
%   author = {{Spring Security Team}},
%   title  = {Spring Security Reference},
%   year   = {2024},
%   url    = {https://docs.spring.io/spring-security/reference}
% }
% @misc{springdatajpa,
%   author = {{Spring Data Team}},
%   title  = {Spring Data JPA Reference},
%   year   = {2024},
%   url    = {https://docs.spring.io/spring-data/jpa/docs/current/reference/html}
% }
% @misc{thymeleaf,
%   author = {{Thymeleaf Team}},
%   title  = {Thymeleaf 3.x Documentation},
%   year   = {2024},
%   url    = {https://www.thymeleaf.org/documentation.html}
% }
% @misc{bootstrap5,
%   author = {{The Bootstrap Authors}},
%   title  = {Bootstrap 5.3 Documentation},
%   year   = {2024},
%   url    = {https://getbootstrap.com/docs/5.3}
% }
% @misc{bootstrapicons,
%   author = {{The Bootstrap Authors}},
%   title  = {Bootstrap Icons 1.11},
%   year   = {2024},
%   url    = {https://icons.getbootstrap.com}
% }
% @misc{chartjs,
%   author = {{Chart.js Contributors}},
%   title  = {Chart.js 4.x Documentation},
%   year   = {2024},
%   url    = {https://www.chartjs.org/docs/latest}
% }
% @misc{qodana,
%   author = {{JetBrains}},
%   title  = {Qodana — Static Code Analysis},
%   year   = {2024},
%   url    = {https://www.jetbrains.com/qodana}
% }
% @book{nielsen1994,
%   author    = {Nielsen, Jakob},
%   title     = {Usability Engineering},
%   publisher = {Morgan Kaufmann},
%   year      = {1994}
% }
% @techreport{wcag21,
%   author      = {{W3C Web Content Accessibility Guidelines Working Group}},
%   title       = {{WCAG} 2.1},
%   institution = {W3C},
%   year        = {2018},
%   url         = {https://www.w3.org/TR/WCAG21}
% }
% @techreport{nist80063b,
%   author      = {Grassi, Paul A. and others},
%   title       = {{NIST SP 800-63B}: Digital Identity Guidelines -- Authentication},
%   institution = {NIST},
%   year        = {2017}
% }
% @article{pierce2004,
%   author  = {Pierce, Jon L.},
%   title   = {The Role of Icons in User Interface Design},
%   journal = {Journal of Usability Studies},
%   year    = {2004}
% }
% @book{gamma1994,
%   author    = {Gamma, Erich and Helm, Richard and Johnson, Ralph and Vlissides, John},
%   title     = {Design Patterns: Elements of Reusable Object-Oriented Software},
%   publisher = {Addison-Wesley},
%   year      = {1994}
% }
.
% ================================================================


% ################################################################
%  CAPÍTULO 1 — INTRODUÇÃO
%  (insira na seção de motivação / objetivos do capítulo 1)
% ################################################################

% ----------------------------------------------------------------
%  Parágrafo de contextualização — inserir após a apresentação do
%  problema de pesquisa
% ----------------------------------------------------------------
A interface de uma aplicação web é o primeiro ponto de contato entre
o usuário e o sistema; uma interface mal projetada pode comprometer a
adoção da ferramenta independentemente da qualidade do seu código de
negócio~\cite{nielsen1994}.  Nesse contexto, o presente trabalho dedica
atenção especial ao desenvolvimento de uma camada de apresentação
profissional, responsiva e acessível, utilizando tecnologias consolidadas
do ecossistema Java — Thymeleaf, Bootstrap 5 e Chart.js — sobre a
plataforma Spring Boot.

% ----------------------------------------------------------------
%  Objetivo específico — adicionar à lista de objetivos específicos
% ----------------------------------------------------------------
\begin{itemize}
  \item Projetar e implementar interfaces administrativas e um portal do
    cliente com padrões visuais modernos, responsividade \textit{mobile-first}
    e navegação intuitiva, adequadas à apresentação acadêmica e ao uso
    em ambiente de produção.
  \item Incorporar visualização de dados (\textit{dashboard} com gráfico de
    \textit{doughnut} e barras de progresso) que permita ao gestor monitorar
    a distribuição de agendamentos por status em tempo real.
\end{itemize}

% ----------------------------------------------------------------
%  Justificativa — inserir no parágrafo de justificativa
% ----------------------------------------------------------------
A escolha por \textit{server-side rendering} com Thymeleaf, em vez de um
\textit{framework} SPA (\textit{Single Page Application}), justifica-se pela
natureza do projeto: um sistema de agendamento com domínio bem definido,
cuja maioria das interações do usuário segue fluxos lineares (cadastro,
listagem, agendamento). Essa abordagem reduz a complexidade de implantação,
elimina a necessidade de uma API REST separada para as telas internas e
facilita a validação de formulários no servidor com o Spring MVC
\textit{BindingResult}.


% ################################################################
%  CAPÍTULO 2 — FUNDAMENTAÇÃO TEÓRICA
% ################################################################

% ----------------------------------------------------------------
%  Seção: Tecnologias de Backend
%  (inserir na subseção de Spring Boot / Spring MVC)
% ----------------------------------------------------------------
\subsection{Spring Boot e Spring MVC}

O Spring Boot~\cite{springboot} é um \textit{framework} que simplifica a
configuração e a implantação de aplicações Java baseadas no ecossistema
Spring.  Ele adota o paradigma de \textit{convention over configuration}:
por meio de dependências \textit{starter} e auto-configurações, o
desenvolvedor pode iniciar uma aplicação web produtiva com poucas linhas de
configuração.

O Spring MVC implementa o padrão arquitetural
\textit{Model-View-Controller}~\cite{gamma1994}: o \texttt{@Controller}
recebe a requisição HTTP, delega a lógica de negócio à camada de serviço e
devolve um nome de \textit{view} ao \texttt{DispatcherServlet}, que aciona
o motor de \textit{templates} Thymeleaf para gerar o HTML final enviado ao
navegador.

\subsection{Spring Security}

O Spring Security~\cite{springsecurity} fornece autenticação e autorização ao
sistema.  No contexto do painel administrativo, a autenticação é realizada
por meio de formulário HTML com sessão gerenciada pelo servidor; o
\texttt{SecurityConfig} configura as rotas protegidas (\texttt{/dashboard},
\texttt{/web/**}) e as rotas públicas (\texttt{/portal/**}).  O portal do
cliente utiliza autenticação manual via \texttt{HttpSession}, com senhas
armazenadas em \textit{hash} BCrypt (\texttt{BCryptPasswordEncoder}).

\subsection{Spring Data JPA e PostgreSQL}

A persistência é gerenciada pelo Spring Data JPA~\cite{springdatajpa},
que abstrai as operações de banco de dados por meio de interfaces
\texttt{JpaRepository}.  As entidades — \texttt{Cliente},
\texttt{Profissional}, \texttt{Servico} e \texttt{Agendamento} — são
mapeadas com anotações JPA (\texttt{@Entity}, \texttt{@ManyToOne},
\texttt{@Enumerated}) e o banco de dados utilizado é o PostgreSQL.

% ----------------------------------------------------------------
%  Seção: Tecnologias de Interface e Frontend
%  (inserir como subseção de Tecnologias Utilizadas)
% ----------------------------------------------------------------
\subsection{Tecnologias de Interface e \textit{Frontend}}

A camada de apresentação foi construída inteiramente sobre tecnologias
consolidadas do ecossistema Java, garantindo compatibilidade nativa com o
Spring Boot, manutenibilidade e facilidade de evolução futura.

\subsubsection{Thymeleaf}

O Thymeleaf~\cite{thymeleaf} é um motor de \textit{templates} para a
plataforma Java integrado nativamente ao Spring MVC.  Toda a geração de
HTML é realizada no servidor (\textit{server-side rendering}), eliminando a
necessidade de APIs REST separadas para as telas administrativas e o portal
do cliente.  Os \textit{templates} herdam de layouts-base por meio dos
mecanismos de fragmento (\texttt{th:replace} / \texttt{th:fragment}),
garantindo consistência visual e evitando duplicação de código.

Um aspecto técnico relevante explorado neste trabalho diz respeito à
interação entre fragmentos Thymeleaf e JavaScript: quando um fragmento é
passado como parâmetro a um \textit{layout} via
\texttt{th:replace="\textasciitilde\{layout :: html(\textasciitilde\{:: conteudo\})\}"},
qualquer diretiva \texttt{th:inline="javascript"} dentro desse fragmento
não é processada de forma confiável~\cite{thymeleaf}.  A solução adotada
foi injetar os dados do controlador como atributos \texttt{data-*} no
elemento HTML (\texttt{th:attr}) e lê-los no JavaScript puro por meio de
\texttt{element.dataset}, eliminando a dependência do processamento
\textit{inline} e mitigando riscos de XSS por interpolação direta.

\subsubsection{Bootstrap 5}

O Bootstrap 5.3.3~\cite{bootstrap5} fornece o sistema de \textit{grid}
responsivo baseado em \textit{Flexbox}, componentes visuais (Navbar, Modal,
Card, Badge, Table, Progress) e utilitários CSS utilizados em todas as
telas.  A versão 5 eliminou a dependência do jQuery, tornando o
\textit{framework} mais leve e compatível com o JavaScript puro utilizado
nos \textit{scripts} de validação e interação.

\subsubsection{Bootstrap Icons}

A biblioteca Bootstrap Icons 1.11.3~\cite{bootstrapicons} é empregada como
sistema de ícones SVG (\textit{icon font}), presente em botões, títulos de
página, itens de menu e \textit{badges} de status.  O uso consistente de
ícones contextuais acelera o reconhecimento de destinos de navegação e
reduz a carga cognitiva do usuário~\cite{pierce2004}.

\subsubsection{Chart.js}

O Chart.js 4.4.3~\cite{chartjs} é incorporado via CDN para renderizar o
gráfico de \textit{doughnut} de distribuição de agendamentos por status no
\textit{dashboard} administrativo.  Os dados são injetados pelo Thymeleaf
por meio de atributos \texttt{data-*} no elemento \texttt{<canvas>}
(Listagem~\ref{lst:chartjs-dataset}) e lidos no JavaScript por meio de
\texttt{canvas.dataset}, garantindo separação clara entre lógica de
servidor e apresentação no cliente.

\begin{lstlisting}[language=HTML, caption={Injeção de dados via atributos \texttt{data-*}
  e inicialização do Chart.js via \texttt{canvas.dataset}.},
  label=lst:chartjs-dataset, basicstyle=\small\ttfamily]
<!-- Template Thymeleaf: injeta dados no canvas -->
<canvas id="graficoAgendamentosStatus"
        th:attr="data-agendado=${totalAgendado},
                 data-confirmado=${totalConfirmado},
                 data-concluido=${totalConcluido},
                 data-cancelado=${totalCancelado}"
        style="max-height:220px;"></canvas>

<!-- JavaScript puro: le os dados sem th:inline -->
<script>
(function () {
    var c = document.getElementById('graficoAgendamentosStatus');
    var agendado   = parseInt(c.dataset.agendado   || '0', 10);
    var confirmado = parseInt(c.dataset.confirmado || '0', 10);
    var concluido  = parseInt(c.dataset.concluido  || '0', 10);
    var cancelado  = parseInt(c.dataset.cancelado  || '0', 10);
    new Chart(c, { type: 'doughnut',
        data: { labels: ['AGENDADO','CONFIRMADO','CONCLUIDO','CANCELADO'],
            datasets: [{ data: [agendado, confirmado, concluido, cancelado],
                backgroundColor: ['#667eea','#28a745','#6c757d','#dc3545'],
                borderWidth: 2, borderColor: '#fff' }] },
        options: { plugins: { legend: { display: false } },
            cutout: '65%', responsive: true,
            maintainAspectRatio: false } });
})();
</script>
\end{lstlisting}

% ----------------------------------------------------------------
%  Seção: Usabilidade e Design de Interface
% ----------------------------------------------------------------
\subsection{Usabilidade e \textit{Design} de Interface}

\subsubsection{Heurísticas de Nielsen}

Nielsen~\cite{nielsen1994} definiu dez heurísticas de usabilidade para o
\textit{design} de interfaces de usuário, amplamente adotadas como
referência na avaliação e no projeto de sistemas interativos.  Dentre as
heurísticas mais diretamente aplicadas neste trabalho destacam-se:
\textit{Consistency and Standards} (uso uniforme de cores, ícones e padrões
de interação), \textit{Recognition rather than recall} (ícones contextuais
nos menus) e \textit{User control and freedom} (botões de cancelar/voltar
em todos os formulários e modais de confirmação para ações destrutivas).

\subsubsection{Responsividade \textit{Mobile-First}}

A abordagem \textit{mobile-first} preconiza que o \textit{layout} base seja
projetado para telas pequenas e progressivamente aprimorado para telas
maiores por meio de \textit{media queries}~\cite{bootstrap5}.  O sistema
emprega o \textit{grid} Bootstrap (\texttt{col-sm-*}, \texttt{col-lg-*})
para adaptar todas as telas ao viewport de 375 px (smartphone) sem overflow
horizontal.


% ################################################################
%  CAPÍTULO 3 (ou 4) — METODOLOGIA
% ################################################################

% ----------------------------------------------------------------
%  Subseção: Implementação da Interface — inserir em Metodologia
% ----------------------------------------------------------------
\subsection{Implementação da Camada de Interface}

O desenvolvimento da camada de interface seguiu um processo iterativo
composto por quatro etapas principais:

\begin{enumerate}
  \item \textbf{Diagnóstico} — inspeção de cada tela existente identificando
    deficiências visuais, HTML malformado e ausência de padrão entre
    componentes.
  \item \textbf{Priorização} — as correções foram ordenadas por impacto no
    usuário: primeiro os defeitos estruturais (HTML inválido, rotas
    quebradas), em seguida melhorias de usabilidade (formulários, tabelas,
    modais) e por último aprimoramentos visuais e adição de funcionalidades
    novas (gráfico de \textit{doughnut}).
  \item \textbf{Implementação} — as alterações foram realizadas arquivo a
    arquivo, com compilação Maven (\texttt{mvn compile}) após cada grupo de
    mudanças para garantir ausência de erros de \textit{template}.
  \item \textbf{Verificação} — cada tela foi verificada segundo o
    \textit{checklist} de testes de interface (Apêndice~A), que cobre os
    fluxos principais e os estados de erro em resolução \textit{desktop}
    (1280 px) e \textit{mobile} (375 px).
\end{enumerate}

\subsection{Controle de Qualidade Estático}

O projeto é analisado automaticamente pelo Qodana~\cite{qodana}, ferramenta
de análise estática da JetBrains integrada ao pipeline CI/CD do repositório
GitHub.  Durante o desenvolvimento foram identificados e corrigidos dois
grupos de avisos:

\begin{itemize}
  \item \textbf{Bulk operation} — substituição de laços \texttt{for} que
    populavam mapas linha a linha por chamadas \texttt{stream().collect
    (Collectors.toMap(...))} e \texttt{stream().collect
    (Collectors.groupingBy(...))}, conforme recomendação de estilo da
    plataforma Java.
  \item \textbf{Nullability / data flow} — substituição de comparador
    lambda com acesso encadeado sem verificação de nulo por
    \texttt{Comparator.comparing(AgendamentoResponse::getDataHora)},
    garantindo que a referência de método é nula-segura em tempo de compilação.
\end{itemize}

\subsection{Estratégia de Ramificação (\textit{Branching})}

O desenvolvimento seguiu o fluxo GitHub Flow: cada conjunto de alterações é
realizado em uma \textit{branch} de funcionalidade (\texttt{feature/}),
submetida como \textit{Pull Request} para revisão automatizada (Qodana) e
mesclada à \textit{branch} principal (\texttt{main}) após aprovação dos
\textit{checks} de CI.


% ################################################################
%  CAPÍTULO DE RESULTADOS E DISCUSSÃO
% ################################################################

\section{Resultados e Discussão}

Esta seção apresenta os resultados obtidos com a implementação da camada de
interface, analisando cada módulo do sistema sob os aspectos funcional,
visual e de usabilidade.

% ......................
\subsection{Painel Administrativo}

\subsubsection{Login}

A tela de login (\texttt{/web/login}, Figura~\ref{fig:tela-login}) passou a
exibir um fundo escuro com gradiente e um \textit{card} centralizado de alta
profundidade visual (\texttt{box-shadow: 0 24px 60px rgba(0,0,0,0.4)}).  O
campo de senha conta com um botão de alternância de visibilidade
(\textit{password toggle}), funcionalidade recomendada pelo NIST SP
800-63B~\cite{nist80063b} para redução de erros de autenticação.

\begin{figure}[H]
  \centering
  \includegraphics[width=0.55\textwidth]{figuras/tela-login.png}
  \caption{Tela de login com gradiente escuro e alternância de senha.}
  \label{fig:tela-login}
\end{figure}

\subsubsection{Navbar e Layout Base}

A barra de navegação (\texttt{layout/base.html}, Figura~\ref{fig:admin-navbar})
foi reconstruída com Bootstrap Navbar responsiva, ícones semânticos
Bootstrap Icons em cada item de menu, destaque dinâmico do item ativo por
JavaScript e colapso em \textit{hamburger menu} em telas menores que 992 px.
Alertas de sucesso e erro são exibidos com ícone e botão de fechamento
(\texttt{data-bs-dismiss="alert"}).

\begin{figure}[H]
  \centering
  \includegraphics[width=0.95\textwidth]{figuras/admin-navbar.png}
  \caption{Navbar administrativa com ícones e responsividade.}
  \label{fig:admin-navbar}
\end{figure}

\subsubsection{\textit{Dashboard}}

O \textit{dashboard} (\texttt{/dashboard}, Figura~\ref{fig:dashboard})
concentra as métricas operacionais mais relevantes em uma única tela:

\begin{itemize}
  \item \textbf{Quatro cartões de métricas} com fundo em gradiente, valor
    numérico em tipografia \textit{bold} de grande porte e efeito
    \textit{hover} de elevação (transform translateY(-4px)).
  \item \textbf{Gráfico de \textit{doughnut}} (Chart.js) com distribuição
    percentual dos agendamentos por status (Figura~\ref{fig:dashboard-chart}).
  \item \textbf{Barras de progresso Bootstrap} como visualização alternativa
    nativa do servidor (independente de JavaScript), exibindo contagem e
    percentual de cada status.
  \item \textbf{Tabela de próximos agendamentos} — até três registros com
    status AGENDADO ou CONFIRMADO, ordenados cronologicamente.
  \item \textbf{Acesso rápido} — três \textit{cards} com contadores de
    clientes, profissionais e serviços, com botões \textit{Gerenciar} e
    \textit{Novo} diretos.
\end{itemize}

\begin{figure}[H]
  \centering
  \includegraphics[width=0.95\textwidth]{figuras/dashboard.png}
  \caption{\textit{Dashboard} com métricas, gráfico e acesso rápido.}
  \label{fig:dashboard}
\end{figure}

\begin{figure}[H]
  \centering
  \includegraphics[width=0.55\textwidth]{figuras/dashboard-chart.png}
  \caption{Gráfico \textit{doughnut} de distribuição de agendamentos por status.}
  \label{fig:dashboard-chart}
\end{figure}

A Tabela~\ref{tab:cores-status} detalha a paleta de cores adotada para os
status, consistente entre o gráfico, as barras de progresso e os
\textit{badges} das tabelas.

\begin{table}[H]
\centering
\begin{tabular}{lll}
\hline
\textbf{Status} & \textbf{Cor}  & \textbf{Código Hexadecimal} \\ \hline
AGENDADO        & Azul-violeta  & \texttt{\#667eea}           \\
CONFIRMADO      & Verde         & \texttt{\#28a745}           \\
CONCLUÍDO       & Cinza         & \texttt{\#6c757d}           \\
CANCELADO       & Vermelho      & \texttt{\#dc3545}           \\ \hline
\end{tabular}
\caption{Paleta de cores dos status de agendamento.}
\label{tab:cores-status}
\end{table}

\paragraph{Desafio técnico — Chart.js em fragmento Thymeleaf.}
A integração do Chart.js com Thymeleaf exigiu a resolução de um problema
não documentado de forma proeminente: a diretiva \texttt{th:inline="javascript"}
não é processada quando o fragmento que a contém é passado como parâmetro a
um \textit{layout} template.  A solução implementada utilizou atributos
\texttt{data-*} no elemento \texttt{<canvas>} (injetados via \texttt{th:attr}
no servidor) e leitura via \texttt{canvas.dataset} no JavaScript do cliente,
conforme descrito na Listagem~\ref{lst:chartjs-dataset}.  Esta abordagem é
igualmente recomendada por questões de segurança: evita interpolação direta
de dados em contexto JavaScript, que poderia introduzir vulnerabilidades XSS.

\subsubsection{Listagens Administrativas}

Todas as telas de listagem seguem um padrão visual unificado
(Figuras~\ref{fig:clientes-lista},~\ref{fig:agendamentos-lista}):
cabeçalho com contador e botão de criação; busca em tempo real via JavaScript
(clientes, profissionais, serviços) ou filtro por formulário GET (agendamentos);
tabela Bootstrap \texttt{table-hover} com botões de ação ícone-apenas;
modal de confirmação para exclusão.

A listagem de agendamentos apresentava HTML malformado na versão anterior —
\texttt{<table>} sem classes Bootstrap, \texttt{</div>} órfãos e referências
a modais inexistentes.  A tela foi reescrita na íntegra, preservando toda a
lógica de transição de status e exclusão.

\begin{figure}[H]
  \centering
  \includegraphics[width=0.95\textwidth]{figuras/clientes-lista.png}
  \caption{Listagem de clientes com busca em tempo real e ações por ícone.}
  \label{fig:clientes-lista}
\end{figure}

\begin{figure}[H]
  \centering
  \includegraphics[width=0.95\textwidth]{figuras/agendamentos-lista.png}
  \caption{Listagem de agendamentos com filtros, badges de status e modais.}
  \label{fig:agendamentos-lista}
\end{figure}

\subsection{Portal do Cliente}

\subsubsection{\textit{Home} do Portal}

A página inicial do portal (\texttt{/portal}, Figura~\ref{fig:portal-home})
foi reformulada com um cartão \textit{hero} em gradiente
(\texttt{linear-gradient(135deg, \#667eea 0\%, \#764ba2 100\%)}), substituindo
o \textit{layout} simples anterior.  Três \textit{cards} de estatísticas
exibem serviços disponíveis, total de profissionais e agendamentos do cliente
autenticado.

\begin{figure}[H]
  \centering
  \includegraphics[width=0.95\textwidth]{figuras/portal-home.png}
  \caption{\textit{Home} do portal com \textit{hero card} e estatísticas.}
  \label{fig:portal-home}
\end{figure}

\subsubsection{\textit{Wizard} de Agendamento}

O fluxo de agendamento em quatro etapas
(Figuras~\ref{fig:wizard-step1},~\ref{fig:wizard-step3},~\ref{fig:wizard-step4})
guia o cliente desde a seleção do serviço até a confirmação com número de
protocolo.  O indicador de progresso visual (\textit{stepper}) comunica em
qual etapa o cliente se encontra e quais já foram concluídas (ícone de
verificação verde).  A Etapa 1 continha dois blocos \texttt{<script>}
duplicados que registravam os mesmos \textit{event listeners}; o bloco
redundante foi removido, eliminando a execução duplicada das funções de
seleção.

\begin{figure}[H]
  \centering
  \includegraphics[width=0.75\textwidth]{figuras/wizard-step1.png}
  \caption{Etapa 1 do \textit{wizard}: seleção de serviço por \textit{card}.}
  \label{fig:wizard-step1}
\end{figure}

\begin{figure}[H]
  \centering
  \includegraphics[width=0.65\textwidth]{figuras/wizard-step3.png}
  \caption{Etapa 3 do \textit{wizard}: revisão e confirmação do agendamento.}
  \label{fig:wizard-step3}
\end{figure}

\begin{figure}[H]
  \centering
  \includegraphics[width=0.65\textwidth]{figuras/wizard-step4.png}
  \caption{Etapa 4: tela de sucesso com protocolo de confirmação.}
  \label{fig:wizard-step4}
\end{figure}

\subsection{Melhorias de Usabilidade — Síntese}

A Tabela~\ref{tab:melhorias} resume as principais melhorias implementadas,
associando cada uma à heurística de Nielsen correspondente.

\begin{table}[H]
\centering
\small
\begin{tabular}{p{5cm} p{4cm} p{4.5cm}}
\hline
\textbf{Melhoria} & \textbf{Heurística} & \textbf{Resultado} \\ \hline
Fundo gradiente + card sombra na tela de login &
  Aesthetic \& minimalist design &
  Foco no formulário, redução de distração \\[4pt]
Ícones nos itens de menu &
  Recognition rather than recall &
  Navegação mais rápida \\[4pt]
\textit{Password toggle} &
  Error prevention &
  Menos erros de digitação de senha \\[4pt]
Modal de confirmação para exclusão &
  User control and freedom &
  Reversibilidade de ações destrutivas \\[4pt]
Gráfico + barras de progresso no dashboard &
  Visibility of system status &
  Distribuição de agendamentos em tempo real \\[4pt]
Busca em tempo real nas listagens &
  Efficiency of use &
  Localização rápida de registros \\[4pt]
Correção HTML malformado (agendamentos) &
  Consistency and standards &
  Comportamento previsível dos componentes \\ \hline
\end{tabular}
\caption{Melhorias de usabilidade e heurísticas de Nielsen correspondentes.}
\label{tab:melhorias}
\end{table}

\subsection{Discussão}

Os resultados demonstram que a adoção de um \textit{design system} baseado
em variáveis CSS globais e nos componentes Bootstrap 5 foi fundamental para
manter a consistência visual entre as mais de quinze telas do sistema, sem
impacto na lógica de negócio.  A separação entre camada de apresentação
(Thymeleaf + Bootstrap) e camada de negócio (serviços e repositórios Spring)
permaneceu preservada em todas as alterações realizadas.

A decisão de oferecer barras de progresso Bootstrap como alternativa nativa
ao gráfico Chart.js mostrou-se acertada sob dois aspectos: primeiro, garante
que o \textit{dashboard} exiba informações de status mesmo em ambientes com
JavaScript desabilitado ou bloqueado; segundo, evidencia os dados numéricos
de forma inequívoca, complementando a visualização proporcional do gráfico.

O desafio técnico da integração Chart.js/Thymeleaf ilustra um aspecto
relevante do desenvolvimento com \textit{server-side rendering}: a
necessidade de compreender os limites do processamento de \textit{templates}
no lado do servidor antes de tomar decisões de integração com bibliotecas
JavaScript.  A solução por atributos \texttt{data-*} é idiomática na web
moderna e agnóstica ao \textit{framework} de \textit{templates}, podendo ser
reaproveitada em projetos futuros que utilizem Freemarker, Pebble ou JSP.


% ################################################################
%  CAPÍTULO — CONCLUSÃO
% ################################################################

\section{Conclusão}

Este trabalho apresentou o desenvolvimento de um sistema de agendamento
construído sobre a plataforma Spring Boot, com ênfase na qualidade da camada
de interface e na experiência do usuário.  O sistema atende ao objetivo
principal de oferecer um painel administrativo completo e um portal do
cliente com fluxo de agendamento guiado, ambos responsivos e visualmente
adequados à apresentação acadêmica e ao ambiente de produção.

\subsection{Contribuições}

As principais contribuições técnicas do trabalho são:

\begin{itemize}
  \item \textbf{Arquitetura de \textit{templates} Thymeleaf} com fragmentos e
    layouts-base reutilizáveis, demonstrando boas práticas de organização de
    \textit{templates} em projetos Spring Boot.
  \item \textbf{Padrão de integração Chart.js/Thymeleaf} via atributos
    \texttt{data-*}, documentando e resolvendo a incompatibilidade entre
    \texttt{th:inline="javascript"} e fragmentos parametrizados, que não
    consta de forma destacada na documentação oficial do Thymeleaf.
  \item \textbf{\textit{Design system} coerente} — paleta, espaçamentos, raios
    de borda e sombras definidos em variáveis CSS (\texttt{:root}), facilmente
    adaptável a outras identidades visuais.
  \item \textbf{Fluxo de agendamento em \textit{wizard}} com indicador de
    progresso, validação por etapa e tela de confirmação com protocolo,
    padrão aplicável a outros domínios de negócio.
  \item \textbf{Pipeline de qualidade estática} — correção de avisos Qodana
    sobre operações de coleção e nulabilidade, demonstrando integração de
    análise estática no fluxo de desenvolvimento.
\end{itemize}

\subsection{Limitações}

O trabalho apresenta as seguintes limitações que podem ser endereçadas em
trabalhos futuros:

\begin{itemize}
  \item \textbf{Testes de interface automatizados} — o \textit{checklist} de
    testes foi realizado manualmente.  A automação com Selenium ou Playwright
    aumentaria a cobertura e a rastreabilidade das verificações.
  \item \textbf{Cobertura de testes de unidade} — os testes automatizados
    existentes cobrem os serviços principais; a camada de \textit{controllers}
    não possui testes de integração com MockMvc.
  \item \textbf{Acessibilidade formal} — embora as interfaces sigam
    recomendações gerais de contraste WCAG 2.1, não foi realizada auditoria
    formal com ferramentas como Axe ou Lighthouse.
  \item \textbf{Internacionalização} — o sistema está fixado em português
    brasileiro; suporte a múltiplos idiomas exigiria a adoção do mecanismo
    \texttt{MessageSource} do Spring.
\end{itemize}

\subsection{Trabalhos Futuros}

Como extensões naturais do sistema, destacam-se:

\begin{itemize}
  \item Implementação de notificações por e-mail e/ou WhatsApp para lembrete
    de agendamentos com proximidade da data.
  \item Exposição de uma API REST com autenticação JWT para integração com
    aplicativos móveis nativos.
  \item Adição de relatórios exportáveis (PDF/CSV) de agendamentos por
    período, profissional ou serviço.
  \item Suporte a múltiplos estabelecimentos (multi-tenant) na mesma
    instalação.
\end{itemize}

Em síntese, o sistema de agendamento desenvolvido neste trabalho cumpre seu
objetivo acadêmico e funcional, entregando uma plataforma operacional com
interfaces de qualidade profissional, arquitetura limpa e pipeline de
qualidade integrado.


% ################################################################
%  REFERÊNCIAS — adicionar ao arquivo .bib do projeto
% ################################################################

%
% @misc{springboot,
%   author = {{Pivotal Software / VMware}},
%   title  = {Spring Boot Reference Documentation},
%   year   = {2024},
%   url    = {https://docs.spring.io/spring-boot/docs/current/reference/html}
% }
% @misc{springsecurity,
%   author = {{Spring Security Team}},
%   title  = {Spring Security Reference},
%   year   = {2024},
%   url    = {https://docs.spring.io/spring-security/reference}
% }
% @misc{springdatajpa,
%   author = {{Spring Data Team}},
%   title  = {Spring Data JPA Reference},
%   year   = {2024},
%   url    = {https://docs.spring.io/spring-data/jpa/docs/current/reference/html}
% }
% @misc{thymeleaf,
%   author = {{Thymeleaf Team}},
%   title  = {Thymeleaf 3.x Documentation},
%   year   = {2024},
%   url    = {https://www.thymeleaf.org/documentation.html}
% }
% @misc{bootstrap5,
%   author = {{The Bootstrap Authors}},
%   title  = {Bootstrap 5.3 Documentation},
%   year   = {2024},
%   url    = {https://getbootstrap.com/docs/5.3}
% }
% @misc{bootstrapicons,
%   author = {{The Bootstrap Authors}},
%   title  = {Bootstrap Icons 1.11},
%   year   = {2024},
%   url    = {https://icons.getbootstrap.com}
% }
% @misc{chartjs,
%   author = {{Chart.js Contributors}},
%   title  = {Chart.js 4.x Documentation},
%   year   = {2024},
%   url    = {https://www.chartjs.org/docs/latest}
% }
% @misc{qodana,
%   author = {{JetBrains}},
%   title  = {Qodana — Static Code Analysis},
%   year   = {2024},
%   url    = {https://www.jetbrains.com/qodana}
% }
% @book{nielsen1994,
%   author    = {Nielsen, Jakob},
%   title     = {Usability Engineering},
%   publisher = {Morgan Kaufmann},
%   year      = {1994}
% }
% @techreport{wcag21,
%   author      = {{W3C Web Content Accessibility Guidelines Working Group}},
%   title       = {{WCAG} 2.1},
%   institution = {W3C},
%   year        = {2018},
%   url         = {https://www.w3.org/TR/WCAG21}
% }
% @techreport{nist80063b,
%   author      = {Grassi, Paul A. and others},
%   title       = {{NIST SP 800-63B}: Digital Identity Guidelines -- Authentication},
%   institution = {NIST},
%   year        = {2017}
% }
% @article{pierce2004,
%   author  = {Pierce, Jon L.},
%   title   = {The Role of Icons in User Interface Design},
%   journal = {Journal of Usability Studies},
%   year    = {2004}
% }
% @book{gamma1994,
%   author    = {Gamma, Erich and Helm, Richard and Johnson, Ralph and Vlissides, John},
%   title     = {Design Patterns: Elements of Reusable Object-Oriented Software},
%   publisher = {Addison-Wesley},
%   year      = {1994}
% }
